% SPDX-License-Identifier: MIT
% Copyright (c) 2026 Evgeny Baulin
% Optimization Methods, Seminar 03: theory handout (English).
% Generated by checks/03. Gradient Descent and Automatic Differentiation.ipynb from web/03/js/content-*.js and content/students.json. Do not edit by hand.
% Compile with: latexmk -xelatex -interaction=nonstopmode -halt-on-error Theory_en.tex
\documentclass[11pt,a4paper]{article}

\usepackage{fontspec}
\setmainfont{cmunrm.otf}[BoldFont=cmunbx.otf, ItalicFont=cmunti.otf, BoldItalicFont=cmunbi.otf]
\setsansfont{cmunss.otf}[BoldFont=cmunsx.otf, ItalicFont=cmunsi.otf, BoldItalicFont=cmunso.otf]
\setmonofont{cmuntt.otf}[BoldFont=cmuntb.otf, ItalicFont=cmunit.otf, BoldItalicFont=cmuntx.otf]
\usepackage{polyglossia}
\setdefaultlanguage{english}

\usepackage[a4paper,left=2.3cm,right=2.3cm,top=2.4cm,bottom=2.4cm,headheight=14pt]{geometry}
\usepackage{amsmath,amssymb}
\usepackage{graphicx}
\usepackage{array,booktabs,longtable}
\usepackage{enumitem}
\usepackage{xcolor}
\usepackage[most]{tcolorbox}
\usepackage{titlesec}
\usepackage{fancyhdr}
\usepackage{needspace}
\usepackage[expansion=false]{microtype}
\usepackage[colorlinks=true,linkcolor=hseblue,urlcolor=hseblue,citecolor=hseblue]{hyperref}

\definecolor{hseblue}{RGB}{16,45,105}
\definecolor{hseaccent}{RGB}{198,40,40}
\definecolor{hsegray}{RGB}{108,114,126}
\definecolor{hselight}{RGB}{240,243,248}
\definecolor{hsewarm}{RGB}{253,244,232}

\hypersetup{pdftitle={Optimization Methods. Seminar 03. Theory handout}, pdfauthor={Evgeny Baulin}}

\titleformat{\section}{\color{hseblue}\Large\bfseries}{\thesection}{0.6em}{}
\titleformat{\subsection}{\color{hseblue}\large\bfseries}{\thesubsection}{0.6em}{}
\titlespacing*{\section}{0pt}{16pt plus 4pt minus 2pt}{8pt plus 2pt}
\titlespacing*{\subsection}{0pt}{12pt plus 3pt minus 2pt}{6pt plus 2pt}
% a heading keeps a few lines under it on its page: a breakable box right after it would otherwise leave it alone
\pretocmd{\section}{\Needspace{7\baselineskip}}{}{\errmessage{cannot patch section}}
\pretocmd{\subsection}{\Needspace{5\baselineskip}}{}{\errmessage{cannot patch subsection}}

\tcbset{handoutbox/.style={enhanced, breakable, sharp corners, boxrule=0pt, left=7pt, right=7pt, top=5pt, bottom=5pt, before skip=8pt plus 2pt, after skip=8pt plus 2pt}}
\newtcolorbox{examplebox}{handoutbox, colback=hselight, colframe=hselight, before upper={\textcolor{hseblue}{\bfseries Example.}\ }}
\newtcolorbox{whybox}{handoutbox, colback=white, colframe=hseblue, leftrule=2.5pt, top=3pt, bottom=3pt, before upper={\textcolor{hseaccent}{\bfseries Why.}\ }}
\newtcolorbox{answerbox}{handoutbox, colback=hsewarm, colframe=hsewarm, before upper={\textcolor{hseaccent}{\bfseries Answer.}\ }}
\newtcolorbox{conventionbox}{handoutbox, colback=hselight, colframe=hseblue, leftrule=2.5pt, before upper={\textcolor{hseblue}{\bfseries Notation.}\par\smallskip}}

\newcommand{\handoutfigure}[3]{\par\medskip\noindent\begin{minipage}{\linewidth}\centering
  \includegraphics[width=#1\linewidth]{figures/#2.pdf}\par\smallskip
  \begin{minipage}{0.94\linewidth}\small\textcolor{hseblue}{\bfseries Figure.}\ #3\end{minipage}\end{minipage}\par\medskip}
\newcommand{\problemhead}[3]{\par\bigskip\noindent{\color{hseblue}\large\bfseries #1.\ #2}\par\nopagebreak
  \noindent{\small\itshape\color{hsegray}#3}\par\nopagebreak\smallskip}
\newtcolorbox{definitionbox}[1]{handoutbox, colback=white, colframe=hseblue, leftrule=2.5pt, before upper={\textcolor{hseblue}{\bfseries #1}\ }}
\newtcolorbox{theorembox}[1]{handoutbox, colback=hselight, colframe=hseblue, leftrule=2.5pt, before upper={\textcolor{hseblue}{\bfseries #1}\ }}
\newtcolorbox{proofbox}{handoutbox, colback=white, colframe=hsegray, leftrule=1.2pt, top=3pt, bottom=3pt, before upper={\textcolor{hsegray}{\bfseries\itshape Proof.}\ }}
\newtcolorbox{notebox}[1]{handoutbox, colback=hselight, colframe=hselight, before upper={\textcolor{hseblue}{\bfseries #1}\par\smallskip}}
\newtcolorbox{trapbox}{handoutbox, colback=hsewarm, colframe=hseaccent, leftrule=2.5pt, before upper={\textcolor{hseaccent}{\bfseries Trap.}\ }}
\newtcolorbox{selfcheckbox}{handoutbox, colback=white, colframe=hseaccent, leftrule=2.5pt, top=3pt, bottom=3pt, before upper={\textcolor{hseaccent}{\bfseries Check yourself.}\ }}
% a display wider than the line is set as a centred paragraph that may break after a relation
\newcommand{\handoutdisplay}[1]{\sbox0{$\displaystyle #1$}\ifdim\wd0>\linewidth
  \par{\centering\binoppenalty=10000\relpenalty=0\relax$\displaystyle #1$\par}\else\[#1\]\fi}
\newcommand{\selfcheckanswer}{\par\smallskip\textcolor{hsegray}{\bfseries Answer.}\ }

\setlist[itemize]{itemsep=3pt,topsep=4pt}
\setlist[enumerate]{itemsep=3pt,topsep=4pt}
\setlength{\parindent}{0pt}
\setlength{\parskip}{5pt plus 1pt}
\setlength{\emergencystretch}{2em}
\setcounter{tocdepth}{2}

\pagestyle{fancy}
\fancyhf{}
\fancyhead[L]{\small\color{hsegray}Optimization Methods}
\fancyhead[R]{\small\color{hsegray}Seminar 03}
\fancyfoot[L]{\small\color{hsegray}Evgeny Baulin}
\fancyfoot[C]{\small\thepage}
\fancyfoot[R]{\small\color{hsegray}MIT License}
\renewcommand{\headrulewidth}{0pt}
\renewcommand{\footrulewidth}{0pt}
\fancypagestyle{plain}{\pagestyle{fancy}}

\begin{document}

\begin{center}
  {\large\color{hsegray}Optimization Methods\par}
  \vspace{4pt}
  {\large\bfseries\color{hseaccent}Seminar 03\par}
  \vspace{8pt}
  {\LARGE\bfseries\color{hseblue}Gradient Descent and Automatic Differentiation\par}
  \vspace{8pt}
  {\Large Theory handout\par}
  \vspace{12pt}
  HSE University, Faculty of Computer Science\par
  Academic year 2026–2027\par
  \vspace{8pt}
  Course lecturer: Andrey Ignatov\par
  Evgeny Baulin\par
  Telegram: \href{https://t.me/tarakan_tuc}{@tarakan\_tuc}\par
\end{center}

\noindent{\color{hseblue}\rule{\linewidth}{0.6pt}}

\tableofcontents

\clearpage


These notes collect the theory of Seminar 3: the definitions for the colloquium, Jensen's inequality, gradient descent with a constant step and with backtracking, convergence rates and reverse-mode automatic differentiation. Every statement comes with its hypotheses and, where the colloquium may ask for it, a complete proof, and the seminar problems appear with full solutions.

\begin{conventionbox}
Vectors are columns, $x^\top y$ is the inner product and $\|\cdot\|$ the Euclidean norm unless a subscript says otherwise. The function $f$ is defined on $\mathbb R^n$, $f^* = \inf_x f(x)$ is the optimal value and $x^*$ a minimizer when one exists. The iterates are $x_0, x_1, \dots$, the step is $\alpha$ or $\alpha_k$, the direction $d$ or $d_k$, and $x^+$ is the next iterate inside a proof. The constants of strong convexity and smoothness are $\mu$ and $L$, and $\kappa = L/\mu$. For symmetric matrices $A \preceq B$ means that $B - A$ is positive semidefinite. In automatic differentiation $\dot v$ is a tangent and $\bar v = \partial z/\partial v$ an adjoint.
\end{conventionbox}

% =====================================================================
\section{Definitions}\label{sec:definitions}

At the colloquium definitions come first, and a definition counts only with its hypotheses. Each block below is the form to reproduce; the short paragraphs between them are consequences, not part of the definitions.

\begin{definitionbox}{Definition: Local and global minimum.}
Let $f: \mathbb R^n \to \mathbb R$. A point $x^*$ is a \textbf{global minimizer} of $f$ if $f(x^*) \le f(x)$ for all $x \in \mathbb R^n$. It is a \textbf{local minimizer} if there is $\varepsilon > 0$ such that $f(x^*) \le f(x)$ for all $x$ with $\|x - x^*\| < \varepsilon$. A minimizer is \textbf{strict} if the inequality is strict for every $x \ne x^*$ in the corresponding set. The value $f(x^*)$ is then a global or local \textbf{minimum}, and $f^* = \inf_x f(x)$ is the optimal value.
\end{definitionbox}

The optimal value may be $-\infty$, and it may be finite without being attained: $f(x) = e^x$ on $\mathbb R$ has $f^* = 0$ and no minimizer.

\begin{definitionbox}{Definition: Stationary point.}
Let $f: \mathbb R^n \to \mathbb R$ be differentiable at $x$. The point $x$ is \textbf{stationary} if $\nabla f(x) = 0$.
\end{definitionbox}

A local minimizer of a differentiable function is stationary: if $\nabla f(x^*) \ne 0$, then $d = -\nabla f(x^*)$ is a descent direction (defined below) and $f$ decreases along it. The converse fails, since a stationary point may be a local maximizer or a saddle point. For a convex differentiable $f$ every stationary point is a global minimizer: put $\nabla f(x^*) = 0$ into the tangent inequality of the next section.

\begin{definitionbox}{Definition: Convex set.}
A set $C \subseteq \mathbb R^n$ is \textbf{convex} if $(1-t)x + ty \in C$ for all $x, y \in C$ and all $t \in [0, 1]$, that is, if it contains the segment between any two of its points.
\end{definitionbox}

\begin{definitionbox}{Definition: Convex and strictly convex function.}
Let $C \subseteq \mathbb R^n$ be a convex set, usually $C = \mathbb R^n$. A function $f: C \to \mathbb R$ is \textbf{convex} if it satisfies Jensen's inequality
\handoutdisplay{
f\big((1-t)x + ty\big) \le (1-t)f(x) + t f(y) \quad \text{for all } x, y \in C,\ t \in [0, 1].
}
 It is \textbf{strictly convex} if the inequality is strict whenever $x \ne y$ and $t \in (0, 1)$. It is \textbf{concave} if $-f$ is convex.
\end{definitionbox}

\begin{definitionbox}{Definition: $\mu$-strong convexity.}
Let $\mu > 0$. A function $f: \mathbb R^n \to \mathbb R$ is $\mu$-\textbf{strongly convex} if $f(x) - \frac{\mu}{2}\|x\|^2$ is convex. Equivalently, Jensen's inequality holds with a quadratic margin:
\handoutdisplay{
f\big((1-t)x + ty\big) \le (1-t)f(x) + t f(y) - \frac{\mu}{2}\,t(1-t)\|x - y\|^2 \quad \text{for all } x, y,\ t \in [0, 1].
}
 If $f$ is differentiable, this is equivalent to the first-order form
\handoutdisplay{
f(y) \ge f(x) + \nabla f(x)^\top (y - x) + \frac{\mu}{2}\|y - x\|^2 \quad \text{for all } x, y,
}
 and if $f$ is twice differentiable, to $\nabla^2 f(x) \succeq \mu I$ for all $x$.
\end{definitionbox}

The forms agree because $(1-t)\|x\|^2 + t\|y\|^2 - \|(1-t)x + ty\|^2 = t(1-t)\|x - y\|^2$ (the identity of the next section), and the first-order form is the tangent inequality for $f - \frac{\mu}{2}\|\cdot\|^2$. A strongly convex function is strictly convex and has exactly one minimizer. Strict convexity is weaker: $f(x) = \log(1 + e^x)$ has $f''(x) = \sigma(x)\big(1 - \sigma(x)\big) > 0$ with $\sigma(x) = 1/(1 + e^{-x})$, but $f''(x) \to 0$ as $|x| \to \infty$, so no $\mu > 0$ works, and this $f$ has no minimizer at all.

\begin{definitionbox}{Definition: $L$-smoothness.}
Let $L > 0$. A differentiable function $f: \mathbb R^n \to \mathbb R$ is $L$-\textbf{smooth} if its gradient is $L$-Lipschitz:
\handoutdisplay{
\|\nabla f(x) - \nabla f(y)\| \le L\|x - y\| \quad \text{for all } x, y.
}
 If $f$ is twice continuously differentiable, this is equivalent to $-LI \preceq \nabla^2 f(x) \preceq LI$ for all $x$. For a convex $f$ the left inequality is automatic, and the condition reads $\nabla^2 f(x) \preceq LI$.
\end{definitionbox}

The lecture uses the two names, $L$-smooth and $L$-Lipschitz gradient, for the same property. The property that the convergence proofs use is the descent lemma of the gradient descent section: $f(y) \le f(x) + \nabla f(x)^\top (y - x) + \frac{L}{2}\|y - x\|^2$.

\begin{definitionbox}{Definition: Condition number.}
Let $f$ be $\mu$-strongly convex and $L$-smooth, with $\mu$ the largest and $L$ the smallest constants for which this holds. The \textbf{condition number} of $f$ is $\kappa = L/\mu$, and $\kappa \ge 1$. For a quadratic $f(x) = \tfrac{1}{2} x^\top Q x - b^\top x$ with symmetric $Q \succ 0$ the best constants are $\mu = \lambda_{\min}(Q)$ and $L = \lambda_{\max}(Q)$, so $\kappa = \lambda_{\max}(Q)/\lambda_{\min}(Q)$.
\end{definitionbox}

The level sets of such a quadratic are ellipsoids, and their longest and shortest semi-axes are in the ratio $\sqrt{\kappa}$. A large $\kappa$ means a long narrow valley.

\begin{definitionbox}{Definition: Descent direction.}
Let $f: \mathbb R^n \to \mathbb R$ be differentiable at $x$. A vector $d \in \mathbb R^n$ is a \textbf{descent direction} at $x$ if $\nabla f(x)^\top d < 0$. Then $f(x + \alpha d) < f(x)$ for all sufficiently small $\alpha > 0$, since $f(x + \alpha d) = f(x) + \alpha\,\nabla f(x)^\top d + o(\alpha)$.
\end{definitionbox}

If $\nabla f(x) \ne 0$, the vector $d = -\nabla f(x)$ is a descent direction: $\nabla f(x)^\top d = -\|\nabla f(x)\|^2 < 0$. The condition is sufficient for local decrease, not necessary: at the maximizer $0$ of $f(x) = -x^2$ every $d \ne 0$ decreases $f$, while $f'(0)\,d = 0$.

\begin{definitionbox}{Definition: Polyak–Łojasiewicz condition.}
Let $f: \mathbb R^n \to \mathbb R$ be differentiable with $f^* = \inf_x f(x) > -\infty$, and let $\mu > 0$. The function $f$ satisfies the \textbf{Polyak–Łojasiewicz (PL) condition} with constant $\mu$ if
\handoutdisplay{
\frac{1}{2}\|\nabla f(x)\|^2 \ge \mu\big(f(x) - f^*\big) \quad \text{for all } x.
}

\end{definitionbox}

Under PL every stationary point is a global minimizer, since $\nabla f(x) = 0$ forces $f(x) = f^*$. A $\mu$-strongly convex differentiable function satisfies PL with the same $\mu$: minimizing the right-hand side of the first-order form over $y$, at $y = x - \nabla f(x)/\mu$, gives $f^* \ge f(x) - \frac{1}{2\mu}\|\nabla f(x)\|^2$. The converse fails, and PL does not even need convexity: $f(x) = x^2 + 3\sin^2 x$ satisfies PL with some $\mu > 0$, while $f''(\pi/2) = 2 + 6\cos\pi < 0$.

% =====================================================================
\section{Jensen's inequality}\label{sec:jensen}

Jensen's inequality is the definition of convexity. The usual hypotheses of the convergence results are the same inequality with a quadratic margin: a lower one for strong convexity, an upper one for smoothness.

\subsection{Reading the inequality}\label{sub:jensen-reading}

Convexity of $f$ on a convex set is Jensen's inequality
\handoutdisplay{
f\big((1-t)x + ty\big) \le (1-t)f(x) + t f(y), \qquad t \in [0, 1].
}
 As $t$ runs from $0$ to $1$, the point $(1-t)x + ty$ runs over the segment from $x$ to $y$, so the left-hand side is the graph of $f$ over the segment. The right-hand side applies the same weights to the values: it is the chord from $(x, f(x))$ to $(y, f(y))$. Convexity says exactly that the chord never dips below the graph. Differentiability plays no role here: $f(x) = |x|$ is convex by the triangle inequality $|(1-t)x + ty| \le (1-t)|x| + t|y|$, yet it is not differentiable at $0$.

The restriction $t \in [0, 1]$ matters. For other $t$ the point $(1-t)x + ty$ still lies on the line through $x$ and $y$, but not between them, and the inequality fails even for $f(x) = x^2$. With $x = 0$, $y = 1$ and $t = 2$ the point is $2$, the left-hand side is $f(2) = 4$, and the right-hand side is $(1-t)f(x) + tf(y) = 2$.

\begin{theorembox}{Proposition: The identity for $x^2$.}
For all real $x$, $y$ and $t$,
\handoutdisplay{
(1-t)x^2 + ty^2 - \big((1-t)x + ty\big)^2 = t(1-t)(x - y)^2.
}

\end{theorembox}

\begin{proofbox}
Expand $\big((1-t)x + ty\big)^2 = (1-t)^2x^2 + 2t(1-t)xy + t^2y^2$ and use $(1-t) - (1-t)^2 = t(1-t)$ and $t - t^2 = t(1-t)$: the difference is $t(1-t)(x^2 - 2xy + y^2) = t(1-t)(x - y)^2$.
\end{proofbox}

For $t \in [0, 1]$ the right-hand side is nonnegative, so $x^2$ is convex, and it vanishes only for $t \in \{0, 1\}$ or $x = y$, so $x^2$ is strictly convex. For $t \notin [0, 1]$ the factor $t(1-t)$ is negative and the inequality fails for every $x \ne y$; in the example above $t(1-t)(x - y)^2 = 2 - 4 < 0$. The same computation with inner products gives, for vectors,
\handoutdisplay{
(1-t)\|x\|^2 + t\|y\|^2 - \|(1-t)x + ty\|^2 = t(1-t)\|x - y\|^2,
}
 and for $f(x) = \tfrac{1}{2} x^\top Q x$ with symmetric $Q$ the gap between chord and graph is $\tfrac{1}{2}\,t(1-t)(x - y)^\top Q (x - y)$.

\subsection{Three equivalent characterizations}\label{sub:jensen-characterizations}

The chord is the definition. Two other tests are equivalent to it, and each needs smoothness that the definition does not.

\begin{center}\small
\begin{tabular}{@{}>{\raggedright\arraybackslash\hspace{0pt}}p{0.25\linewidth}>{\raggedright\arraybackslash\hspace{0pt}}p{0.22\linewidth}>{\raggedright\arraybackslash\hspace{0pt}}p{0.42\linewidth}@{}}
\toprule
Characterization & Needs & Meaning \\
\midrule
Chord (Jensen): $f\big((1-t)x + ty\big) \le (1-t)f(x) + tf(y)$ & a convex domain, no derivatives & the graph lies on or below every chord \\
Tangent: $f(y) \ge f(x) + \nabla f(x)^\top (y - x)$ & $f$ differentiable on an open convex set & the graph lies on or above every tangent plane; the linear model is a global lower bound \\
Curvature: $\nabla^2 f(x) \succeq 0$ & $f$ twice differentiable on an open convex set & $d^\top \nabla^2 f(x)\, d \ge 0$ for every $d$: the graph curves upward along every line \\
\bottomrule
\end{tabular}
\end{center}

\begin{theorembox}{Theorem: Characterizations of convexity.}
Let $C \subseteq \mathbb R^n$ be open and convex. If $f$ is differentiable on $C$, it is convex if and only if $f(y) \ge f(x) + \nabla f(x)^\top (y - x)$ for all $x, y \in C$. If $f$ is twice differentiable on $C$, it is convex if and only if $\nabla^2 f(x) \succeq 0$ for all $x \in C$.
\end{theorembox}

\begin{proofbox}
From the chord to the tangent: for $t \in (0, 1]$ the chord inequality gives $f\big(x + t(y - x)\big) - f(x) \le t\big(f(y) - f(x)\big)$; divide by $t$ and let $t \to 0$ to get $\nabla f(x)^\top (y - x) \le f(y) - f(x)$. From the tangent to the chord: for $z = (1-t)x + ty$ the tangent inequality at $z$ gives $f(x) \ge f(z) + \nabla f(z)^\top (x - z)$ and $f(y) \ge f(z) + \nabla f(z)^\top (y - z)$; multiply by $1-t$ and $t$ and add, and the gradient terms cancel because $(1-t)(x - z) + t(y - z) = 0$. Curvature: $f$ is convex if and only if every restriction $\varphi(s) = f(x + sd)$ to a line is convex, a twice differentiable function of one variable is convex if and only if $\varphi'' \ge 0$, and $\varphi''(s) = d^\top \nabla^2 f(x + sd)\, d$.
\end{proofbox}

For strict convexity only one direction survives: $\nabla^2 f(x) \succ 0$ for all $x$ implies strict convexity, but $f(x) = |x|^3$ is strictly convex while $f''(x) = 6|x|$ vanishes at $0$.

\subsection{Finite and expectation forms}\label{sub:jensen-forms}

\begin{theorembox}{Theorem: Jensen's inequality for several points.}
Let $f$ be convex on a convex set $C$, let $x_1, \dots, x_N \in C$, and let $t_1, \dots, t_N \ge 0$ with $\sum_i t_i = 1$. Then
\handoutdisplay{
f\Big(\sum_{i=1}^N t_i x_i\Big) \le \sum_{i=1}^N t_i f(x_i).
}

\end{theorembox}

\begin{proofbox}
Induction on $N$. For $N = 1$ both sides are $f(x_1)$, and $N = 2$ is the definition. Let $N \ge 3$ and assume the claim for $N - 1$ points. If $t_N = 1$, all other weights vanish and there is nothing to prove. Otherwise
\handoutdisplay{
\sum_{i=1}^N t_i x_i = (1 - t_N)\,y + t_N x_N, \qquad y = \sum_{i=1}^{N-1} \frac{t_i}{1 - t_N}\,x_i,
}
 and $y$ is a convex combination of $N - 1$ points of $C$: its weights are nonnegative and sum to $1$, so $y \in C$. The definition and then the induction hypothesis give
\handoutdisplay{
f\Big(\sum_{i=1}^N t_i x_i\Big) \le (1 - t_N) f(y) + t_N f(x_N) \le \sum_{i=1}^{N-1} t_i f(x_i) + t_N f(x_N).
}

\end{proofbox}

\begin{theorembox}{Theorem: Jensen's inequality for expectations.}
Let $f: \mathbb R^n \to \mathbb R$ be convex and let $X$ be a random vector in $\mathbb R^n$ with $\mathbb E\|X\| < \infty$ and $\mathbb E|f(X)| < \infty$. Then
\handoutdisplay{
f(\mathbb E X) \le \mathbb E f(X).
}

\end{theorembox}

\begin{proofbox}
Put $m = \mathbb E X$. If $f$ is differentiable, the tangent inequality at $m$ reads $f(X) \ge f(m) + \nabla f(m)^\top (X - m)$; take expectations and use $\mathbb E(X - m) = 0$. Without differentiability the same argument works with a subgradient at $m$, a vector $g$ with $f(y) \ge f(m) + g^\top (y - m)$ for all $y$, which exists at every point for a convex $f: \mathbb R^n \to \mathbb R$. A random vector that takes the value $x_i$ with probability $t_i$ gives back the finite form.
\end{proofbox}

\subsection{Concave functions and the AM–GM inequality}\label{sub:jensen-concave}

\begin{trapbox}
For a concave $g$ every inequality above reverses: $g\big(\sum_i t_i x_i\big) \ge \sum_i t_i g(x_i)$ and $g(\mathbb E X) \ge \mathbb E g(X)$. The logarithm is concave, so $\mathbb E \log X \le \log \mathbb E X$ for a positive random variable $X$, and not the other way round.
\end{trapbox}

\begin{theorembox}{Proposition: AM–GM for two numbers.}
For $a, b > 0$,
\handoutdisplay{
\sqrt{ab} \le \frac{a + b}{2},
}
 with equality if and only if $a = b$.
\end{theorembox}

\begin{proofbox}
The function $-\log$ is strictly convex on the convex set $(0, \infty)$, since $(-\log x)'' = 1/x^2 > 0$. The chord inequality at $t = \tfrac{1}{2}$ gives
\handoutdisplay{
-\log\frac{a + b}{2} \le \frac{1}{2}\big(-\log a - \log b\big) = -\log\sqrt{ab},
}
 strictly unless $a = b$. Multiply by $-1$ and apply the increasing function $\exp$.
\end{proofbox}

\begin{examplebox}
$a = 1$, $b = 9$: $\sqrt{ab} = 3 \le 5 = \frac{a + b}{2}$. With the finite form the same argument gives the weighted version $\prod_i a_i^{t_i} \le \sum_i t_i a_i$ for $a_i > 0$, $t_i \ge 0$, $\sum_i t_i = 1$.
\end{examplebox}

\subsection{A counterexample}\label{sub:jensen-counterexample}

\begin{examplebox}
Take $f(x) = x^3$ on $\mathbb R$ with $x = -2$, $y = 1$ and $t = 0.5$. The point is $(1-t)x + ty = -0.5$, the left-hand side is $f(-0.5) = -0.125$, and the right-hand side is $(1-t)x^3 + ty^3 = -3.5$. Since $-0.125 > -3.5$, the inequality fails, and $x^3$ is not convex on $\mathbb R$.
\end{examplebox}

One pair and one $t$ disprove convexity; no number of successful checks proves it. On $[0, \infty)$ the function $x^3$ is convex, since $f''(x) = 6x \ge 0$ there, so a random check that draws 10000 pairs from $[0, 2]$ finds 0 violations and says nothing about negative arguments. Convexity is proved with the definition or with one of the characterizations above.

\subsection{Quadratic margins on both sides}\label{sub:jensen-margins}

In the language of the chord, $\mu$ and $L$ bound the gap between the chord and the graph from below and from above. If $f$ is $\mu$-strongly convex and $L$-smooth, then for all $x, y$ and $t \in [0, 1]$
\handoutdisplay{
\frac{\mu}{2}\,t(1-t)\|x - y\|^2 \le (1-t)f(x) + tf(y) - f\big((1-t)x + ty\big) \le \frac{L}{2}\,t(1-t)\|x - y\|^2,
}
 and in first-order form
\handoutdisplay{
\frac{\mu}{2}\|y - x\|^2 \le f(y) - f(x) - \nabla f(x)^\top (y - x) \le \frac{L}{2}\|y - x\|^2.
}
 The left inequalities are the definition of strong convexity. The right one in first-order form is the descent lemma, and its chord version follows by writing the descent lemma at $z = (1-t)x + ty$ for the points $x$ and $y$ and adding with weights $1-t$ and $t$. For a convex differentiable $f$ each right inequality is in fact equivalent to $L$-smoothness.

For $f(x) = \tfrac{1}{2} x^\top Q x$ the gap between chord and graph equals $\tfrac{1}{2}\,t(1-t)(x - y)^\top Q (x - y)$. It lies between the two bounds with $\mu = \lambda_{\min}(Q)$ and $L = \lambda_{\max}(Q)$ and touches each of them when $x - y$ is a corresponding eigenvector. The ratio $\kappa = L/\mu$ says how far apart the two parabolas are, and the next sections show that it controls the number of gradient descent iterations.

\begin{selfcheckbox}
Apply Jensen's inequality to $f(x) = x^2$ and a random variable $X$ with finite variance. What does it say, and what is the gap between the two sides?
\selfcheckanswer It says $(\mathbb E X)^2 \le \mathbb E X^2$. The gap is $\mathbb E X^2 - (\mathbb E X)^2 = \operatorname{Var} X \ge 0$, and it vanishes only when $X$ is constant with probability one. If $X = x$ with probability $1-t$ and $X = y$ with probability $t$, then $\operatorname{Var} X = t(1-t)(x - y)^2$: the identity for $x^2$ is this variance.
\end{selfcheckbox}

% =====================================================================
\section{Gradient descent}\label{sec:gd}

An iterative method starts from $x_0$ and repeats
\handoutdisplay{
d_k = \mathcal D(x_k, f, \nabla f, \dots), \qquad \alpha_k = \mathcal S(x_k, d_k), \qquad x_{k+1} = x_k + \alpha_k d_k.
}
 Three design questions define a method: the direction $d_k$, the step $\alpha_k$ and the rule for stopping. Gradient descent takes $d_k = -\nabla f(x_k)$. This section treats a constant step; backtracking and stopping come in the next one.

\begin{definitionbox}{Definition: Gradient descent.}
Let $f: \mathbb R^n \to \mathbb R$ be differentiable and $\alpha > 0$. Gradient descent with constant step $\alpha$ from $x_0$ is the sequence
\handoutdisplay{
x_{k+1} = x_k - \alpha\nabla f(x_k), \qquad k = 0, 1, 2, \dots
}
 Each iteration costs one gradient evaluation.
\end{definitionbox}

\subsection{Why the negative gradient}\label{sub:gd-steepest}

\begin{theorembox}{Theorem: Steepest ascent in the Euclidean norm.}
Let $f$ be differentiable at $x$ with $\nabla f(x) \ne 0$. Then
\handoutdisplay{
\arg\max_{\|d\|_2 \le 1}\ \lim_{t \downarrow 0} \frac{f(x + td) - f(x)}{t} = \frac{\nabla f(x)}{\|\nabla f(x)\|_2},
}
 and the maximal rate of increase is $\|\nabla f(x)\|_2$.
\end{theorembox}

\begin{proofbox}
Differentiability gives $f(x + td) = f(x) + t\,\nabla f(x)^\top d + o(t)$, so the limit equals $\nabla f(x)^\top d$. By the Cauchy–Schwarz inequality
\handoutdisplay{
\nabla f(x)^\top d \le \|\nabla f(x)\|_2\,\|d\|_2 \le \|\nabla f(x)\|_2,
}
 with equality for $d = \nabla f(x)/\|\nabla f(x)\|_2$. Equality in Cauchy–Schwarz forces $d$ to be a nonnegative multiple of $\nabla f(x)$, and equality in the second step forces $\|d\|_2 = 1$, so this maximizer is unique.
\end{proofbox}

Replacing $d$ by $-d$, the steepest descent direction among unit vectors is $-\nabla f(x)/\|\nabla f(x)\|_2$, with rate $-\|\nabla f(x)\|_2$.

The word \emph{steepest} refers to the Euclidean unit ball. For a general norm the steepest unit direction solves $\max_{\|d\| \le 1} \nabla f(x)^\top d = \|\nabla f(x)\|_*$, where $\|\cdot\|_*$ is the dual norm, and the answer changes with the norm. For $\|d\|_M = \sqrt{d^\top M d}$ with $M \succ 0$ substitute $u = M^{1/2} d$: the problem becomes $\max_{\|u\|_2 \le 1} \big(M^{-1/2}\nabla f(x)\big)^\top u$, the theorem gives $u \propto M^{-1/2}\nabla f(x)$, and so $d \propto M^{-1}\nabla f(x)$. The steepest descent direction in this norm points along $-M^{-1}\nabla f(x)$, the idea behind preconditioning and the natural gradient. For $\|\cdot\|_\infty$ one steepest descent direction is $-\operatorname{sign}\nabla f(x)$, taken componentwise.

\subsection{The descent lemma and a safe step}\label{sub:gd-lemma}

\begin{theorembox}{Lemma: Descent lemma.}
Let $f: \mathbb R^n \to \mathbb R$ be differentiable with an $L$-Lipschitz gradient. Then for all $x, y$
\handoutdisplay{
f(y) \le f(x) + \nabla f(x)^\top (y - x) + \frac{L}{2}\|y - x\|^2.
}

\end{theorembox}

\begin{proofbox}
Put $r(t) = f\big(x + t(y - x)\big)$ for $t \in [0, 1]$. It is continuously differentiable with $r'(t) = \nabla f\big(x + t(y - x)\big)^\top (y - x)$, so
\handoutdisplay{
f(y) - f(x) = \int_0^1 r'(t)\,dt = \nabla f(x)^\top (y - x) + \int_0^1 \big(\nabla f(x + t(y - x)) - \nabla f(x)\big)^\top (y - x)\,dt.
}
 By Cauchy–Schwarz and the Lipschitz condition the integrand is at most $\|\nabla f(x + t(y - x)) - \nabla f(x)\|\,\|y - x\| \le Lt\|y - x\|^2$, and $\int_0^1 Lt\,dt = \frac{L}{2}$.
\end{proofbox}

The same estimate with absolute values gives $\big|f(y) - f(x) - \nabla f(x)^\top (y - x)\big| \le \frac{L}{2}\|y - x\|^2$: an $L$-smooth function stays between two parabolas around its tangent plane. Gradient descent needs only the upper one. The statement is global: one constant $L$ controls the error of the linear model at every point and for every step length.

\begin{theorembox}{Corollary: One-step decrease.}
Let $f$ be $L$-smooth, $\alpha > 0$ and $x^+ = x - \alpha\nabla f(x)$. Then
\handoutdisplay{
f(x^+) \le f(x) - \alpha\Big(1 - \frac{L\alpha}{2}\Big)\|\nabla f(x)\|^2.
}

\end{theorembox}

\begin{proofbox}
Put $y = x^+$ in the descent lemma. Since $y - x = -\alpha\nabla f(x)$,
\handoutdisplay{
f(x^+) \le f(x) - \alpha\|\nabla f(x)\|^2 + \frac{L\alpha^2}{2}\|\nabla f(x)\|^2.
}

\end{proofbox}

Two constants follow. For $0 < \alpha < 2/L$ the factor $1 - L\alpha/2$ is positive, so every step from a non-stationary point strictly decreases $f$. The guaranteed decrease $\alpha(1 - L\alpha/2)$ is a concave parabola in $\alpha$ with its maximum at $\alpha = 1/L$, where
\handoutdisplay{
f(x^+) \le f(x) - \frac{1}{2L}\|\nabla f(x)\|^2,
}
 the best guarantee this bound can give. Every $\alpha \le 1/L$ has $1 - L\alpha/2 \ge \tfrac{1}{2}$, hence $f(x^+) \le f(x) - \frac{\alpha}{2}\|\nabla f(x)\|^2$.

\begin{trapbox}
The value $2/L$ is where the guarantee of strict decrease ends, not a proof that larger steps diverge: at $\alpha = 2/L$ the bound only rules out an increase, for $\alpha > 2/L$ it allows one, and a particular function may still decrease. For a quadratic, $2/L$ is the exact boundary, as the next subsection shows.
\end{trapbox}

\subsection{Quadratics and the exact stability condition}\label{sub:gd-quadratic}

Take $f(x) = \tfrac{1}{2} x^\top Q x - b^\top x$ with symmetric $Q \succ 0$. Then $\nabla f(x) = Qx - b$, the unique minimizer is $x^* = Q^{-1}b$, and $\nabla f(x) = Q(x - x^*)$. For the error $e_k = x_k - x^*$ one step of gradient descent gives
\handoutdisplay{
e_{k+1} = e_k - \alpha Q e_k = (I - \alpha Q)\,e_k.
}
 Diagonalize $Q = U\Lambda U^\top$ with $U$ orthogonal and $\Lambda = \operatorname{diag}(\lambda_1, \dots, \lambda_n)$, and pass to the eigen-coordinates $z_k = U^\top e_k$. The coordinates decouple:
\handoutdisplay{
z_{k+1,i} = (1 - \alpha\lambda_i)\,z_{k,i}, \qquad z_{k,i} = (1 - \alpha\lambda_i)^k z_{0,i},
}
 and $f(x_k) - f^* = \tfrac{1}{2} e_k^\top Q e_k = \tfrac{1}{2}\sum_i \lambda_i z_{k,i}^2$.

\begin{theorembox}{Theorem: Stability on a quadratic.}
Let $f(x) = \tfrac{1}{2} x^\top Q x - b^\top x$ with symmetric $Q \succ 0$. Gradient descent with a constant step $\alpha$ converges to $x^*$ from every starting point if and only if
\handoutdisplay{
0 < \alpha < \frac{2}{\lambda_{\max}(Q)}.
}

\end{theorembox}

\begin{proofbox}
By the recursion, $z_{k,i} \to 0$ for every $z_{0,i}$ if and only if $|1 - \alpha\lambda_i| < 1$, and $\|x_k - x^*\| = \|z_k\|$ because $U$ is orthogonal. So $x_k \to x^*$ from every $x_0$ if and only if $|1 - \alpha\lambda_i| < 1$ for every $i$. Since $\lambda_i > 0$, this means $0 < \alpha\lambda_i < 2$ for all $i$, and the binding eigenvalue is the largest.
\end{proofbox}

At the boundary $\alpha = 2/\lambda_{\max}$ the factor of the stiffest coordinate is $-1$: that coordinate flips sign with constant magnitude forever, while the others still shrink. Beyond the boundary it grows geometrically. Between $1/\lambda_{\max}$ and $2/\lambda_{\max}$ its factor is negative but less than $1$ in magnitude, and it converges with alternating signs. Here $L = \lambda_{\max}(Q)$, so the safe range $0 < \alpha < 2/L$ from the descent lemma is exactly the stability range: for quadratics the lemma is sharp.

\begin{examplebox}
In A1, $Q = \operatorname{diag}(1,\ 10)$, so $L = 10$ and the boundary is $2/L = 0.2$. At $\alpha = 0.2$ the factors are $0.8$ and $-1$. From $x_0 = (2,\ 2)$ the stiff coordinate alternates between $2$ and $-2$, and
\handoutdisplay{
f(x_k) = 20 + 2 \cdot 0.64^k:
}
 the values decrease towards $20$ and never go below it, because the stiff coordinate alone contributes $\tfrac{1}{2} \cdot 10 \cdot 2^2 = 20$.
\end{examplebox}

\handoutfigure{0.9}{fig_stability}{Values $f(x_k)$, log scale, for $f(x) = \tfrac{1}{2} x^\top Q x$ with $Q = \operatorname{diag}(1,\ 10)$, $x_0 = (2,\ 2)$ and the four steps of A1. Below the boundary $2/L = 0.2$ the values fall geometrically, at $\alpha = 2/L$ they level off at $20$, and above it they grow.}

\subsection{Conditioning and the optimal constant step}\label{sub:gd-conditioning}

For the quadratic, let $\mu = \lambda_{\min}(Q)$ and $L = \lambda_{\max}(Q)$. Per step the norm of the error is multiplied by at most
\handoutdisplay{
\rho(\alpha) = \max_i |1 - \alpha\lambda_i| = \max\big(|1 - \alpha\mu|,\ |1 - \alpha L|\big),
}
 because $|1 - \alpha\lambda|$ is convex in $\lambda$ and its maximum over $[\mu, L]$ sits at an endpoint. With $\|e_k\| = \|z_k\|$ this gives
\handoutdisplay{
\|x_k - x^*\| \le \rho(\alpha)^k\,\|x_0 - x^*\|, \qquad f(x_k) - f^* \le \rho(\alpha)^{2k}\big(f(x_0) - f^*\big).
}


At $\alpha = 1/L$ the factors $1 - \lambda_i/L$ lie in $[0, 1 - 1/\kappa]$, so $\rho(1/L) = 1 - 1/\kappa$. Since $\ln\frac{1}{1 - 1/\kappa} \ge \frac{1}{\kappa}$, the error is below $\varepsilon\|x_0 - x^*\|$ once $k \ge \kappa\ln(1/\varepsilon)$: the iteration count grows in proportion to $\kappa$. Geometrically, a large $\kappa$ is a long narrow valley; the gradient points mostly across it, and a step short enough for the steep direction makes little progress along the flat one. On a quadratic the gap therefore shrinks by at least $(1 - \mu/L)^2$ per step at $\alpha = 1/L$. This is sharper than the general guarantee $f(x_k) - f^* \le (1 - \mu/L)^k\big(f(x_0) - f^*\big)$, which holds at $\alpha = 1/L$ for every $\mu$-strongly convex $L$-smooth function and is proved in the section on rates. A2 and A3 ask for that general guarantee, so their factor is $1 - \mu/L$, not $(1 - \mu/L)^2$.

\begin{theorembox}{Proposition: Optimal constant step.}
For a quadratic with $\mu = \lambda_{\min}(Q) > 0$ and $L = \lambda_{\max}(Q)$, the factor $\rho(\alpha)$ is minimized over $\alpha > 0$ by
\handoutdisplay{
\alpha^* = \frac{2}{L + \mu}, \qquad \rho(\alpha^*) = \frac{L - \mu}{L + \mu} = \frac{\kappa - 1}{\kappa + 1}.
}

\end{theorembox}

\begin{proofbox}
For $\alpha > 0$, $|1 - \alpha\lambda| = \max(1 - \alpha\lambda,\ \alpha\lambda - 1)$; the first term is largest at $\lambda = \mu$ and the second at $\lambda = L$. So $\rho(\alpha) = \max(1 - \alpha\mu,\ \alpha L - 1)$, the maximum of a decreasing and an increasing linear function of $\alpha$. It is smallest where the two are equal: $1 - \alpha\mu = \alpha L - 1$ gives $\alpha^* = 2/(L + \mu)$ and $\rho(\alpha^*) = 1 - \frac{2\mu}{L + \mu} = \frac{L - \mu}{L + \mu}$. For $\alpha < \alpha^*$ the first function is larger than at $\alpha^*$, and for $\alpha > \alpha^*$ the second is, so no other step does as well. Dividing by $\mu$ gives $(\kappa - 1)/(\kappa + 1)$.
\end{proofbox}

At $\alpha^*$ the two extreme factors are $1 - \alpha^*\mu = \rho(\alpha^*)$ and $1 - \alpha^* L = -\rho(\alpha^*)$, equal in magnitude. In A3, $\mu = 1$ and $L = 50$ give $\alpha^* = \tfrac{2}{51} \approx 0.039216$, and $\kappa = 50$ gives $\rho(\alpha^*) = \tfrac{49}{51} \approx 0.960784$, against $1 - 1/\kappa = 0.98$ at $\alpha = 1/L$. When $Q$ has only the eigenvalues $\mu$ and $L$, as in every two-dimensional example of the seminar, each coordinate shrinks by exactly $\rho(\alpha^*)$ in magnitude, and $f(x_k) - f^* = \rho(\alpha^*)^{2k}\big(f(x_0) - f^*\big)$ holds with equality. The same step and factor bound $\|x_k - x^*\|$ for every $\mu$-strongly convex $L$-smooth function, with a longer proof.

\handoutfigure{0.95}{fig_conditioning}{Contours of $f(x) = \tfrac{1}{2} x^\top Q x$ and the iterates of gradient descent with $\alpha = 1/L$ from $x_0 = (1,\ 1)$. Left: $Q = \operatorname{diag}(1,\ 2)$, $\kappa = 2$. Right: $Q = \operatorname{diag}(1,\ 50)$, $\kappa = 50$. The flat coordinate shrinks by $1 - 1/\kappa$ per step, $0.5$ on the left and $0.98$ on the right, and reaching $f(x_k) \le 10^{-6}\,f(x_0)$ takes 10 iterations on the left and 245 on the right.}

\subsection{Seminar problems}\label{sub:gd-problems}

The problems of block A with complete solutions.

\problemhead{A1}{Four steps on one quadratic}{On the board}
\textbf{Statement.} Minimize $f(x) = \tfrac{1}{2} x^\top Q x$ with $Q = \operatorname{diag}(1,\ 10)$ by gradient descent with a constant step $\alpha$ from $x_0 = (2,\ 2)$. \textbf{(a)} Find the stability boundary: the value $\alpha_{\max}$ such that the method converges from every start if and only if $0 < \alpha < \alpha_{\max}$. \textbf{(b)} Classify the steps $\alpha = 0.05$, $0.19$, $0.2$ and $0.21$ through the per-coordinate recurrence $z_{k+1,i} = (1 - \alpha\lambda_i)\,z_{k,i}$: does the method converge, and how?

\textbf{Solution.}
\begingroup\sloppy
\begin{enumerate}[label=\textbf{\arabic*.}, leftmargin=*]
  \item Here $x^* = 0$, $f^* = 0$, and $Q$ is already diagonal, so the eigen-coordinates are the coordinates: $z_k = x_k$, $\mu = \lambda_1 = 1$, $L = \lambda_2 = 10$, $\kappa = 10$. The step $x_{k+1} = x_k - \alpha Q x_k$ splits into $z_{k+1,i} = (1 - \alpha\lambda_i)\,z_{k,i}$, and $f(x_k) = \tfrac{1}{2}\big(\lambda_1 z_{k,1}^2 + \lambda_2 z_{k,2}^2\big)$, so $f(x_0) = 22$.
  \item Convergence from every start needs $|1 - \alpha\lambda_i| < 1$ for both $i$, that is $0 < \alpha < 2/\lambda_i$. The flat coordinate allows $\alpha < 2/\mu = 2$, the stiff one only $\alpha < 2/L = 0.2$. The stiff direction sets the boundary: $0 < \alpha < 0.2$.
  \item The factors of the four steps:
\handoutdisplay{
\begin{array}{c|cc} \alpha & 1 - \alpha\lambda_1 & 1 - \alpha\lambda_2 \\ \hline 0.05 & 0.95 & 0.5 \\ 0.19 & 0.81 & -0.9 \\ 0.2 & 0.8 & -1 \\ 0.21 & 0.79 & -1.1 \end{array}
}

  \item $\alpha = 0.05$: both factors lie in $(0, 1)$, so both coordinates shrink without changing sign. The stiff one halves at every step, the flat one loses only a twentieth. After ten steps $f = 0.7170$, and almost all of it sits in the flat coordinate. Converges, slowly.
  \item $\alpha = 0.19$: the stiff factor $-0.9$ is negative, so $z_{k,2}$ changes sign at every step, but its magnitude still shrinks. Converges with sign flips. Early on the stiff coordinate dominates: it holds $20$ of $f(x_0) = 22$ and shrinks by $|-0.9|$ here against $0.5$ at $\alpha = 0.05$, so this step is behind after ten iterations, $f = 2.4611$ against $0.7170$. In the long run the pace is set by the larger factor in magnitude, $|-0.9|$ here against $0.95$, so it is far ahead after fifty, $f = 5.31 \cdot 10^{-4}$ against $1.18 \cdot 10^{-2}$.
  \item $\alpha = 0.2 = 2/L$: the stiff factor is exactly $-1$. The stiff coordinate flips between $+2$ and $-2$ forever and alone contributes $\tfrac{1}{2} \cdot 10 \cdot 2^2 = 20$ to $f$ at every step. The flat coordinate still shrinks by $0.8$, so $f(x_k) = 20 + 2 \cdot 0.64^k$: it decreases towards $20$ and never goes below it ($20.0231$ after ten steps). Neither convergence to $x^*$ nor divergence.
  \item $\alpha = 0.21$: $|-1.1| > 1$, so $|z_{k,2}|$ grows at every step while its sign alternates, and $f(x_k) \to \infty$: $134.5679$ after ten steps, $2.76 \cdot 10^{5}$ after fifty. The flat coordinate converges, which does not help. Diverges.
  \item The boundary comes from the largest eigenvalue, the speed from the smallest: for every $\alpha < 2/L$ the flat factor is $1 - \alpha\mu > 1 - 2/\kappa$. A large $\kappa$ forces a step that barely moves the flat coordinate.
\end{enumerate}
\endgroup
\begin{answerbox}
Stable for $0 < \alpha < 2/L = 0.2$. $\alpha = 0.05$: converges, slowly in the flat coordinate; $0.19$: converges with sign flips in the stiff coordinate; $0.2$: the stiff coordinate flips between $\pm2$ forever and $f(x_k)$ decreases towards $20$, never below; $0.21$: diverges.
\end{answerbox}

\problemhead{A2}{Guaranteed and observed iterations}{On your own}
\textbf{Statement.} Run gradient descent on $f(x) = \tfrac{1}{2} x^\top Q x$ with $Q = \operatorname{diag}(1,\ 2)$, step $\alpha = 1/L$, from $x_0 = (1,\ 1)$. Report \textbf{(a)} $\kappa$; \textbf{(b)} the factor $q$ in the lecture's guarantee $f(x_k) - f^* \le q^k\,\big(f(x_0) - f^*\big)$; \textbf{(c)} the smallest $k$ for which this bound gives $f(x_k) - f^* \le 10^{-6}\,\big(f(x_0) - f^*\big)$; \textbf{(d)} the smallest $k$ for which the iterates actually satisfy it.

\textbf{Solution.}
\begingroup\sloppy
\begin{enumerate}[label=\textbf{\arabic*.}, leftmargin=*]
  \item $L = \lambda_{\max} = 2$, $\mu = \lambda_{\min} = 1$, so $\kappa = 2$ and $\alpha = 1/L = 0.5$.
  \item The linear rate at $\alpha = 1/L$ is $f(x_k) - f^* \le (1 - \mu/L)^k\,\big(f(x_0) - f^*\big)$, so $q = 1 - \mu/L = 0.5$.
  \item The bound reaches $10^{-6}$ when $0.5^k \le 10^{-6}$, that is $k \ge \dfrac{\ln 10^{-6}}{\ln 0.5} = 19.93$. Round up: $20$ iterations.
  \item Now the run itself, with $f^* = 0$. The factors $1 - \alpha\lambda_i$ are $0.5$ for $\lambda = 1$ and $0$ for $\lambda = 2$: the stiff coordinate is gone after the first step.
  \item So $f(x_0) = \tfrac{1}{2}\big(1 + 2\big) = 1.5$ and $f(x_k) = \tfrac{1}{2} \cdot 0.5^{2k}$ for $k \ge 1$. From then on the gap shrinks by $0.5^2$ per iteration, not by $0.5$.
  \item $f(x_k) \le 10^{-6}\,f(x_0)$ means $0.5^{2k} \le 2 \cdot 1.5 \cdot 10^{-6}$, that is $k \ge \dfrac{\ln\big(2 \cdot 1.5 \cdot 10^{-6}\big)}{2 \ln 0.5} = 9.17$. The run needs $10$ iterations.
  \item The bound is weaker because its proof uses only the descent inequality and $\|\nabla f\|^2 \ge 2\mu\,(f - f^*)$, which hold for every function with these $\mu$ and $L$. On a quadratic at $\alpha = 1/L$ each coordinate term of the gap shrinks by $(1 - \lambda_i/L)^2 \le (1 - \mu/L)^2$.
\end{enumerate}
\endgroup
\begin{answerbox}
$\kappa = 2$, $q = 0.5$, the bound allows $20$ iterations, the run needs $10$.
\end{answerbox}

\problemhead{A3}{An ill-conditioned bowl and the optimal step}{On your own}
\textbf{Statement.} Same setup as A2 with $Q = \operatorname{diag}(1,\ 50)$ and $x_0 = (1,\ 1)$. \textbf{(a)–(d)} Report the same four quantities at $\alpha = 1/L$. \textbf{(e)} Give the optimal constant step $\alpha^* = 2/(L + \mu)$ as a number. \textbf{(f)} With $\alpha^*$ every coordinate of the error $x_k - x^*$ is multiplied at each step by a factor of magnitude at most $\rho = (\kappa - 1)/(\kappa + 1)$. How many iterations does the resulting bound on the gap $f(x_k) - f^*$ need for the same reduction by $10^{-6}$?

\textbf{Solution.}
\begingroup\sloppy
\begin{enumerate}[label=\textbf{\arabic*.}, leftmargin=*]
  \item $L = 50$, $\mu = 1$, $\kappa = 50$, $\alpha = 1/L = 0.02$ and $q = 1 - \mu/L = 0.98$.
  \item Bound: $k \ge \dfrac{\ln 10^{-6}}{\ln 0.98} = 683.84$, so $684$ iterations.
  \item The factors at $\alpha = 1/L$ are $1 - \mu/L = 0.98$ and $1 - L/L = 0$. The stiff coordinate holds $\tfrac{1}{2} \cdot 50$ out of $f(x_0) = 25.5$ and is gone after one step; then $f(x_k) = \tfrac{1}{2} \cdot 0.98^{2k}$ for $k \ge 1$.
  \item $f(x_k) \le 10^{-6}\,f(x_0)$ means $0.98^{2k} \le 2 \cdot 25.5 \cdot 10^{-6}$, that is $k \ge \dfrac{\ln\big(2 \cdot 25.5 \cdot 10^{-6}\big)}{2 \ln 0.98} = 244.6$. The run needs $245$ iterations against the bound's $684$.
  \item Optimal step: $\alpha^* = \dfrac{2}{L + \mu} = \tfrac{2}{51} \approx 0.039216$. The factors are $1 - \alpha^*\mu = \tfrac{49}{51}$ and $1 - \alpha^* L = -\tfrac{49}{51}$: equal magnitudes, $\rho = \dfrac{\kappa - 1}{\kappa + 1} = \tfrac{49}{51} \approx 0.960784$.
  \item The gap is quadratic in the coordinates, so $f(x_k) - f^* \le \rho^{2k}\,\big(f(x_0) - f^*\big)$. The bound needs $k \ge \dfrac{\ln 10^{-6}}{2 \ln \rho} = 172.67$, so $173$ iterations. With $\rho^k$ instead of $\rho^{2k}$ one gets $346$, the count for the distance $\|x_k - x^*\|$.
  \item Here both coordinates shrink by exactly $\rho$ in magnitude, so $f(x_k) = \rho^{2k} f(x_0)$ with equality, and the run also needs $173$ iterations. With $\alpha^*$ the bound and the run coincide; at $\alpha = 1/L$ they differ almost threefold.
\end{enumerate}
\endgroup
\begin{answerbox}
$\kappa = 50$, $q = 0.98$, the bound allows $684$ iterations, the run needs $245$; $\alpha^* = \tfrac{2}{51} \approx 0.039216$, and the bound at $\alpha^*$ allows $173$, exactly what the run needs.
\end{answerbox}

% =====================================================================
\section{Line search and stopping}\label{sec:linesearch}

A constant step needs $L$, which is rarely known, and a step that is safe where $f$ curves most is too short elsewhere. Backtracking chooses the step at every iteration from values of $f$ along the ray.

\subsection{Armijo backtracking}\label{sub:linesearch-armijo}

\begin{definitionbox}{Definition: Armijo backtracking.}
Fix an initial step $\alpha_0 > 0$, a shrinkage factor $\beta \in (0, 1)$ and a constant $c \in (0, 1)$, and let $d$ be a descent direction at $x$. Start with $\alpha = \alpha_0$ and repeat $\alpha \leftarrow \beta\alpha$ until the sufficient decrease (Armijo) condition
\handoutdisplay{
f(x + \alpha d) \le f(x) + c\,\alpha\,\nabla f(x)^\top d
}
 holds. The accepted step is the first $\alpha = \beta^j\alpha_0$, $j = 0, 1, 2, \dots$, that passes. For gradient descent $d = -\nabla f(x)$ and the right-hand side is $f(x) - c\,\alpha\|\nabla f(x)\|^2$. Typical values are $\beta = \tfrac{1}{2}$ and $c = 10^{-4}$.
\end{definitionbox}

Along the ray put $\varphi(\alpha) = f(x + \alpha d)$, so $\varphi(0) = f(x)$ and $\varphi'(0) = \nabla f(x)^\top d < 0$. The Armijo line $\varphi(0) + c\alpha\varphi'(0)$ is the tangent line with its slope reduced by the factor $c$, so for $\alpha > 0$ it lies above the tangent. Since $\varphi(\alpha) = \varphi(0) + \alpha\varphi'(0) + o(\alpha)$, short steps pass; a long step fails when the curvature bends $\varphi$ back above the line. The test rejects steps whose decrease is a negligible fraction of what the slope promises.

\begin{theorembox}{Theorem: Termination of backtracking.}
Let $f$ be $L$-smooth, $\nabla f(x) \ne 0$, $d = -\nabla f(x)$, $\alpha_0 > 0$ and $\beta, c \in (0, 1)$. The Armijo condition holds for every $\alpha \in \big(0, 2(1-c)/L\big]$. Consequently backtracking stops after finitely many reductions, and the accepted step satisfies
\handoutdisplay{
\alpha \ge \min\Big(\alpha_0,\ \frac{2\beta(1-c)}{L}\Big).
}

\end{theorembox}

\begin{proofbox}
By the one-step decrease bound, $f\big(x - \alpha\nabla f(x)\big) \le f(x) - \alpha\big(1 - \frac{L\alpha}{2}\big)\|\nabla f(x)\|^2$. If $\alpha \le 2(1-c)/L$, then $1 - L\alpha/2 \ge c$ and the right-hand side is at most $f(x) - c\alpha\|\nabla f(x)\|^2$: the Armijo condition holds. The trial steps $\beta^j\alpha_0$ tend to $0$, so a finite $j$ brings them into $\big(0, 2(1-c)/L\big]$, and the loop stops there at the latest. If $\alpha_0$ passes, the accepted step is $\alpha_0$. Otherwise the accepted step $\alpha$ follows a rejected trial $\alpha/\beta$, every rejected trial exceeds $2(1-c)/L$, and so $\alpha > 2\beta(1-c)/L$.
\end{proofbox}

The proof also bounds the work: the number of reductions is at most the smallest $j \ge 0$ with $\beta^j\alpha_0 \le 2(1-c)/L$, that is $\big\lceil \log_{1/\beta}\big(\alpha_0 L/(2(1-c))\big) \big\rceil$ when $\alpha_0 > 2(1-c)/L$. The loop may stop earlier, because along a particular direction the condition can hold beyond $2(1-c)/L$. On a quadratic the range is exact: with $g = \nabla f(x)$,
\handoutdisplay{
f(x - \alpha g) = f(x) - \alpha\,g^\top g + \frac{\alpha^2}{2}\,g^\top Q g,
}
 so the Armijo condition holds if and only if $\alpha \le 2(1-c)\,\dfrac{g^\top g}{g^\top Q g}$, and this is at least $2(1-c)/L$ because $g^\top Q g \le L\,g^\top g$.

\begin{examplebox}
In B1, $Q = \operatorname{diag}(1,\ 10)$, $x_0 = (2,\ 2)$, $L = 10$ and $c = 10^{-4}$. The theorem guarantees acceptance for $\alpha \le 2(1-c)/L = 0.19998$, and along $d = -\nabla f(x_0)$ the exact range is $\alpha \le 0.20178$. Halving from $\alpha_0 = 1$, the first trial inside both ranges is $0.125$, reached after 3 halvings; it lies between $1/L = 0.1$ and $2/L = 0.2$.
\end{examplebox}

Without $L$-smoothness backtracking still terminates for any descent direction $d$ of a differentiable $f$, because $f(x + \alpha d) - f(x) - c\alpha\nabla f(x)^\top d = (1 - c)\,\alpha\,\nabla f(x)^\top d + o(\alpha)$ is negative for small $\alpha$. What is lost is the lower bound on the accepted step.

\begin{trapbox}
If $d$ is not a descent direction, short trial steps fail the test and the loop may never end. An implementation checks $\nabla f(x)^\top d < 0$ and caps the number of reductions.
\end{trapbox}

\subsection{Stopping rules}\label{sub:linesearch-stopping}

The lecture uses three scaled tests with tolerances $\varepsilon_g$, $\varepsilon_x$ and $\varepsilon_f$:
\handoutdisplay{
\begin{gathered}
\|\nabla f(x_k)\| \le \varepsilon_g \max\big(1, |f(x_k)|\big), \\[6pt]
\|x_{k+1} - x_k\| \le \varepsilon_x \max\big(1, \|x_k\|\big), \\[6pt]
|f(x_{k+1}) - f(x_k)| \le \varepsilon_f \max\big(1, |f(x_k)|\big).
\end{gathered}
}
 The factor $\max(1, \cdot)$ makes a test relative when the scale is large and absolute when it is below $1$, so a quantity near zero is not asked for an impossible relative accuracy.

In code the tolerance is often split into an absolute and a relative part, $\text{atol} + \text{rtol}\cdot\max(1, s)$, where $s$ is the scale, $|f(x_k)|$ or $\|x_k\|$. B3 uses this form with $\text{atol} = 10^{-8}$ and $\text{rtol} = 10^{-6}$.

The tests measure different things. The change tests watch the motion of the method: for gradient descent $\|x_{k+1} - x_k\| = \alpha_k\|\nabla f(x_k)\|$ and $f(x_k) - f(x_{k+1}) \approx \alpha_k\|\nabla f(x_k)\|^2$, and both are small whenever the step is small, however far the point is from a solution. Small changes can mean convergence, a tiny step, poor scaling or floating-point stagnation. The gradient test looks at the point itself. A small gradient is a first-order certificate of stationarity; under the PL condition it bounds the error in $f$, and for a $\mu$-strongly convex $f$ also the distance to $x^*$:
\handoutdisplay{
f(x) - f^* \le \frac{1}{2\mu}\|\nabla f(x)\|^2, \qquad \|x - x^*\| \le \frac{1}{\mu}\|\nabla f(x)\|.
}
 The first is PL rearranged. For the second, add the first-order form of strong convexity at $x$ and at $x^*$ to get $\big(\nabla f(x) - \nabla f(x^*)\big)^\top (x - x^*) \ge \mu\|x - x^*\|^2$, then use $\nabla f(x^*) = 0$ and Cauchy–Schwarz.

\begin{notebox}{Combined rule}
Stop with success when the gradient test holds, with the tolerance scaled by $\max(1, |f(x_k)|)$ as in the lecture. Treat the change tests only as a stagnation signal: if they hold while the gradient test fails, stop and report that the method stalled, not that it converged. Always add a cap on iterations and function evaluations, and stop when $f$ or $\nabla f$ is not finite.
\end{notebox}

\subsection{A change-based test that stops too early}\label{sub:linesearch-failure}

Problem B3 builds the failure on purpose. Take $f(x) = \tfrac{1}{2} x^\top Q x$ with $Q = \operatorname{diag}(1,\ 50)$, so that $x^* = 0$ and $f^* = 0$, start at $x_0 = (1,\ 1)$ with $f(x_0) = 25.5$, and make one gradient step with the prescribed step $\alpha = 10^{-8}$. It gives $x_1 = (0.99999999,\ 0.9999995)$ and $f(x_1) = 25.49997499$. The three tests with $\text{atol} = 10^{-8}$ and $\text{rtol} = 10^{-6}$ compare as follows.

\begin{center}\small
\begin{tabular}{@{}>{\raggedright\arraybackslash\hspace{0pt}}p{0.24\linewidth}rll@{}}
\toprule
Test & Value & Threshold & Fires \\
\midrule
function change $|f(x_1) - f(x_0)|$ & $2.501 \cdot 10^{-5}$ & $\text{atol} + \text{rtol}\max(1, |f(x_0)|) = 2.551 \cdot 10^{-5}$ & yes \\
step $\|x_1 - x_0\|$ & $5.001 \cdot 10^{-7}$ & $\text{atol} + \text{rtol}\max(1, \|x_0\|) = 1.424 \cdot 10^{-6}$ & yes \\
gradient $\|\nabla f(x_1)\|$ & $50.01$ & $\text{atol} + \text{rtol}\max(1, |f(x_1)|) = 2.551 \cdot 10^{-5}$ & no \\
\bottomrule
\end{tabular}
\end{center}

Both change tests fire after one step, although $f$ has dropped by less than a millionth of its value and the point has hardly moved towards $x^*$. Only the gradient test refuses to stop. The change tests fire because their left-hand sides are proportional to the step: $\|x_1 - x_0\| = \alpha\|\nabla f(x_0)\|$ and $|f(x_1) - f(x_0)| \approx \alpha\|\nabla f(x_0)\|^2$ with $\|\nabla f(x_0)\| = 50.01$, and a tiny step makes both tiny. The gradient test has no factor $\alpha$: it measures the gradient at $x_1$ itself, and a tiny step leaves it almost as large as at $x_0$, $\|\nabla f(x_1)\| \approx \|\nabla f(x_0)\|$.

\subsection{Seminar problems}\label{sub:linesearch-problems}

The problems of block B with complete solutions.

\problemhead{B1}{Armijo backtracking by hand}{On the board}
\textbf{Statement.} Let $f(x) = \tfrac{1}{2} x^\top Qx$ with $Q = \operatorname{diag}(1, 10)$ and $x_0 = (2,\ 2)$. Run Armijo backtracking along $d = -\nabla f(x_0)$ with $\alpha_0 = 1$, $\beta = 0.5$ and $c = 10^{-4}$: try $\alpha = \alpha_0$ first and set $\alpha \leftarrow \beta\alpha$ while $f(x_0 + \alpha d) > f(x_0) + c\alpha\nabla f(x_0)^\top d$. How many halvings happen before the test passes, which step is accepted, and how does it compare with $1/L$ and $2/L$?

\textbf{Solution.}
\begingroup\sloppy
\begin{enumerate}[label=\textbf{\arabic*.}, leftmargin=*]
  \item Data at $x_0$: $f(x_0) = 22$, $\nabla f(x_0) = Qx_0 = (2,\ 20)$ and $\nabla f(x_0)^\top d = -\|\nabla f(x_0)\|^2 = -404$. Here $L = 10$, $1/L = 0.1$ and $2/L = 0.2$.
  \item Along the ray each coordinate gets its own factor: $x_0 + \alpha d = \big(2(1 - \alpha),\ 2(1 - 10\alpha)\big)$. The test reads $f(x_0 + \alpha d) \le 22 - 10^{-4}\cdot404\,\alpha$.
  \item Trials, halving from $\alpha_0$:
\handoutdisplay{
\begin{array}{r|r|r|c} \alpha & f(x_0 + \alpha d) & f(x_0) + c\alpha\nabla f(x_0)^\top d & \text{passes} \\ \hline 1 & 1620 & 21.9596 & \text{no} \\ 0.5 & 320.5 & 21.9798 & \text{no} \\ 0.25 & 46.125 & 21.9899 & \text{no} \\ 0.125 & 2.78125 & 21.99495 & \text{yes} \end{array}
}

  \item The first trial that passes is $\alpha = 0.125$, after $3$ halvings, and the loop stops there. It lies between $1/L = 0.1$ and $2/L = 0.2$, inside the stable range of block A.
  \item Why the loop must stop: the descent lemma gives $f(x_0 + \alpha d) \le f(x_0) - \alpha(1 - L\alpha/2)\|\nabla f(x_0)\|^2$, so the test holds for every $\alpha \le 2(1-c)/L = 0.19998$. Every rejected trial is above this value, so the accepted one is at least $\beta$ times it.
  \item On a quadratic the test can be solved exactly. With $g = \nabla f(x_0)$, $f(x_0 - \alpha g) = f(x_0) - \alpha\|g\|^2 + \tfrac{1}{2}\alpha^2 g^\top Qg$, so the test holds for $\alpha \le 2(1-c)\|g\|^2/(g^\top Qg) = 0.20178$, a little above $0.19998$ because the curvature along $g$ is below $L$. The minimizer along the ray is $\|g\|^2/(g^\top Qg) = 0.1009$; Armijo does not look for it.
\end{enumerate}
\endgroup
\begin{answerbox}
$3$ halvings, accepted $\alpha = 0.125$, between $1/L = 0.1$ and $2/L = 0.2$; every $\alpha \le 2(1-c)/L = 0.19998$ passes.
\end{answerbox}

\problemhead{B2}{Backtracking on a stiffer quadratic}{On your own}
\textbf{Statement.} Let $f(x) = \tfrac{1}{2} x^\top Qx$ with $Q = \operatorname{diag}(1, 50)$ and $x_0 = (2,\ 2)$. Run the same backtracking as in B1 along $d = -\nabla f(x_0)$ with $\alpha_0 = 1$, $\beta = 0.5$ and $c = 10^{-4}$: try $\alpha = \alpha_0$ first and set $\alpha \leftarrow \beta\alpha$ while $f(x_0 + \alpha d) > f(x_0) + c\alpha\nabla f(x_0)^\top d$. How many halvings happen, which step is accepted, and where does it sit relative to $1/L$ and $2/L$?

\textbf{Solution.}
\begingroup\sloppy
\begin{enumerate}[label=\textbf{\arabic*.}, leftmargin=*]
  \item Data: $f(x_0) = 102$, $\nabla f(x_0) = (2,\ 100)$, $\|\nabla f(x_0)\|^2 = 10004$; $L = 50$, $1/L = 0.02$, $2/L = 0.04$. With $x_0 + \alpha d = \big(2(1 - \alpha),\ 2(1 - 50\alpha)\big)$ the test is $f(x_0 + \alpha d) \le 102 - 10^{-4}\cdot10004\,\alpha$.
  \item Trials, halving from $\alpha_0$:
\handoutdisplay{
\begin{array}{r|r|r|c} \alpha & f(x_0 + \alpha d) & f(x_0) + c\alpha\nabla f(x_0)^\top d & \text{passes} \\ \hline 1 & 240100 & 100.9996 & \text{no} \\ 0.5 & 57600.5 & 101.4998 & \text{no} \\ 0.25 & 13226.125 & 101.7499 & \text{no} \\ 0.125 & 2757.78125 & 101.87495 & \text{no} \\ 0.0625 & 453.3203125 & 101.937475 & \text{no} \\ 0.03125 & 33.517578125 & 101.9687375 & \text{yes} \end{array}
}

  \item The first pass is $\alpha = 0.03125$ after $5$ halvings. It lies between $1/L = 0.02$ and $2/L = 0.04$, as in B1.
  \item Guarantee: every $\alpha \le 2(1-c)/L = 0.039996$ passes the test, and this bound is just below $2/L$. The halvings cross it between $0.0625$ and $0.03125$, and the accepted step is at least $\beta$ times it.
  \item Exact threshold on this quadratic: $2(1-c)\|g\|^2/(g^\top Qg) = 0.040012$ with $g = \nabla f(x_0)$, slightly above $2/L$. Armijo checks the decrease of one step, not stability: a trial just above $2/L$ would pass here, although as a constant step it diverges.
  \item Compared with B1, $L$ went from $10$ to $50$ and backtracking needed $5$ halvings instead of $3$. In both problems the accepted step lands between $1/L$ and $2/L$, and backtracking never used $L$.
\end{enumerate}
\endgroup
\begin{answerbox}
$5$ halvings, accepted $\alpha = 0.03125$, between $1/L = 0.02$ and $2/L = 0.04$ and below the guarantee $2(1-c)/L = 0.039996$.
\end{answerbox}

\problemhead{B3}{Three stopping tests after a tiny step}{On your own}
\textbf{Statement.} Let $f(x) = \tfrac{1}{2} x^\top Qx$ with $Q = \operatorname{diag}(1, 50)$ and $x_0 = (1,\ 1)$. Make one gradient step $x_1 = x_0 - \alpha\nabla f(x_0)$ with the prescribed step $\alpha = 10^{-8}$. With $\mathrm{atol} = 10^{-8}$, $\mathrm{rtol} = 10^{-6}$ and Euclidean norms, evaluate the three scaled tests
\handoutdisplay{
\begin{aligned} \text{gradient:}\quad & \|\nabla f(x_1)\| \le \mathrm{atol} + \mathrm{rtol}\cdot\max(1, |f(x_1)|), \\ \text{step:}\quad & \|x_1 - x_0\| \le \mathrm{atol} + \mathrm{rtol}\cdot\max(1, \|x_0\|), \\ \text{function:}\quad & |f(x_1) - f(x_0)| \le \mathrm{atol} + \mathrm{rtol}\cdot\max(1, |f(x_0)|). \end{aligned}
}
 A test fires when its inequality holds. Report $|f(x_1) - f(x_0)|$, $\|\nabla f(x_1)\|$ and the number of tests that fire. Is $x_1$ close to the minimizer? Then state a combined stopping rule and say which quantity it is scaled by.

\textbf{Solution.}
\begingroup\sloppy
\begin{enumerate}[label=\textbf{\arabic*.}, leftmargin=*]
  \item One step: $\nabla f(x_0) = Qx_0 = (1, 50)$, $\|\nabla f(x_0)\| = 50.01$, and $x_1 = x_0 - \alpha\nabla f(x_0) = (0.99999999,\ 0.9999995)$.
  \item Function test: $f(x_0) = 25.5$ and $f(x_1) = 25.49997499$, so $|f(x_1) - f(x_0)| = 2.501 \cdot 10^{-5}$, which matches the first-order value $\alpha\|\nabla f(x_0)\|^2$. The threshold is $\mathrm{atol} + \mathrm{rtol}\cdot25.5 = 2.551 \cdot 10^{-5}$: the test fires.
  \item Step test: $\|x_1 - x_0\| = \alpha\|\nabla f(x_0)\| = 5.001 \cdot 10^{-7}$ against $\mathrm{atol} + \mathrm{rtol}\cdot\max(1, \|x_0\|) = \mathrm{atol} + \mathrm{rtol}\sqrt2 = 1.424 \cdot 10^{-6}$: it fires too.
  \item Gradient test: $\|\nabla f(x_1)\| = 50.01$, almost the same as at $x_0$, against $\mathrm{atol} + \mathrm{rtol}\cdot\max(1, |f(x_1)|) = 2.551 \cdot 10^{-5}$. It does not fire, so $2$ tests out of $3$ fire.
  \item The minimizer is $x^* = 0$ with $f^* = 0$, while $\|x_1 - x^*\| \approx \sqrt2$ and $f(x_1) \approx f(x_0)$. The change tests fired because both changes shrink with $\alpha$ (to first order they are proportional to it): they measure the step, not the distance to $x^*$.
  \item Combined rule: report convergence only when $\|\nabla f(x_k)\| \le \mathrm{atol} + \mathrm{rtol}\cdot\max(1, |f(x_k)|)$, a test scaled by $\max(1, |f(x_k)|)$. If only the change tests fire, stop with a stagnation warning, not a solution; always cap the number of iterations and stop on non-finite values.
\end{enumerate}
\endgroup
\begin{answerbox}
$|f(x_1) - f(x_0)| = 2.501 \cdot 10^{-5}$ and $\|x_1 - x_0\| = 5.001 \cdot 10^{-7}$ are below their thresholds, $\|\nabla f(x_1)\| = 50.01$ is not: $2$ tests fire far from $x^* = 0$. Stop on the gradient test scaled by $\max(1, |f(x_k)|)$ and use the change tests only as a stagnation signal, with an iteration cap.
\end{answerbox}

% =====================================================================
\section{Rates}\label{sec:rates}

The lecture proves three guarantees for gradient descent with the step $\alpha = 1/L$: a gap of order $1/k$ for convex functions, and a gap that shrinks by the factor $1 - \mu/L$ at every step for strongly convex functions and, more generally, under the PL condition. All three proofs start from the descent step of the section on gradient descent: for an $L$-smooth $f$ and $x_{k+1} = x_k - \frac1L \nabla f(x_k)$ the descent lemma gives
\handoutdisplay{
f(x_{k+1}) \le f(x_k) - \frac{1}{2L}\|\nabla f(x_k)\|^2.
}
 In particular $f(x_k)$ never increases. Throughout, $f^* = \inf f$, and $x^*$ is a minimizer when one exists.

\subsection{The convex case}\label{sub:rates-convex}

\begin{theorembox}{Theorem: Rate for convex functions.}
Let $f$ be convex and $L$-smooth, and let $x^*$ be a minimizer, $f^* = f(x^*)$. Gradient descent with $\alpha = 1/L$ satisfies
\handoutdisplay{
f(x_k) - f^* \le \frac{L\|x_0 - x^*\|^2}{2k} \qquad \text{for every } k \ge 1.
}

\end{theorembox}

\begin{proofbox}
Write $g_k = \nabla f(x_k)$. Convexity puts the tangent plane at $x_k$ below the graph, so $f^* \ge f(x_k) + g_k^\top(x^* - x_k)$, that is, $f(x_k) \le f^* + g_k^\top(x_k - x^*)$. Substituting this into the descent step gives
\handoutdisplay{
f(x_{k+1}) - f^* \le g_k^\top(x_k - x^*) - \frac{1}{2L}\|g_k\|^2.
}
 The right-hand side is a difference of squared distances. Since $x_{k+1} - x^* = (x_k - x^*) - \frac1L g_k$, expanding the square gives the three-point identity
\handoutdisplay{
\|x_{k+1} - x^*\|^2 = \|x_k - x^*\|^2 - \frac2L\, g_k^\top(x_k - x^*) + \frac{1}{L^2}\|g_k\|^2,
}
 and multiplying it by $L/2$ turns the previous inequality into the key inequality
\handoutdisplay{
f(x_{k+1}) - f^* \le \frac L2\Big(\|x_k - x^*\|^2 - \|x_{k+1} - x^*\|^2\Big).
}
 Summing over the steps $0, 1, \dots, k-1$, the right-hand side telescopes:
\handoutdisplay{
\sum_{j=1}^{k}\big(f(x_j) - f^*\big) \le \frac L2\Big(\|x_0 - x^*\|^2 - \|x_k - x^*\|^2\Big) \le \frac L2\|x_0 - x^*\|^2.
}
 Finally, $f(x_j)$ does not increase, so each of the $k$ terms on the left is at least $f(x_k) - f^*$. Hence $k\big(f(x_k) - f^*\big) \le \frac L2\|x_0 - x^*\|^2$.
\end{proofbox}

The key inequality has a second consequence. Its left-hand side is nonnegative, so $\|x_{k+1} - x^*\| \le \|x_k - x^*\|$: the iterates never move away from any minimizer.

\begin{whybox}
Each step pays for its gap with a decrease of the squared distance to $x^*$, and the whole budget is $\|x_0 - x^*\|^2$. So the gaps are small on average over the first $k$ steps, and since they only decrease, the last one is at most the average.
\end{whybox}

The rate is \textbf{sublinear}: halving the bound takes twice as many steps, and a gap below $\varepsilon$ is guaranteed only after $L\|x_0 - x^*\|^2/(2\varepsilon)$ steps.

\begin{selfcheckbox}
The theorem holds for every minimizer $x^*$. If the minimizers form a whole set, which of them gives the best bound?
\selfcheckanswer The one closest to $x_0$. All minimizers have the same value $f^*$, and the proof uses nothing else about $x^*$, so $\|x_0 - x^*\|$ can be replaced by the distance from $x_0$ to the set of minimizers.
\end{selfcheckbox}

\subsection{Strong convexity}\label{sub:rates-strong}

A differentiable $f$ is $\mu$-\textbf{strongly convex}, $\mu > 0$, if
\handoutdisplay{
f(y) \ge f(x) + \nabla f(x)^\top (y - x) + \frac\mu2\|y - x\|^2 \qquad \text{for all } x, y.
}
 For an $L$-smooth $f$ the descent lemma adds the upper bound, which sandwiches $f(y) - f(x) - \nabla f(x)^\top(y - x)$ between $\frac\mu2\|y - x\|^2$ and $\frac L2\|y - x\|^2$, so $\mu \le L$ and $\kappa = L/\mu \ge 1$.

\begin{theorembox}{Lemma: Strong convexity implies PL.}
If $f$ is $\mu$-strongly convex, then for every $x$
\handoutdisplay{
\frac{1}{2}\|\nabla f(x)\|^2 \ge \mu\big(f(x) - f^*\big).
}

\end{theorembox}

\begin{proofbox}
Fix $x$ and write $g = \nabla f(x)$. The right-hand side of the definition is a quadratic function of $y$,
\handoutdisplay{
q(y) = f(x) + g^\top (y - x) + \frac\mu2\|y - x\|^2.
}
 Its gradient $g + \mu(y - x)$ vanishes at $y = x - g/\mu$, so its smallest value is
\handoutdisplay{
q\Big(x - \frac g\mu\Big) = f(x) - \frac1\mu\|g\|^2 + \frac{1}{2\mu}\|g\|^2 = f(x) - \frac{1}{2\mu}\|g\|^2.
}
 Hence $f(y) \ge q(y) \ge f(x) - \frac{1}{2\mu}\|g\|^2$ for every $y$. Taking the infimum over $y$ gives $f^* \ge f(x) - \frac{1}{2\mu}\|g\|^2$, and multiplying by $\mu$ and rearranging gives the claim.
\end{proofbox}

\begin{theorembox}{Theorem: Linear rate under strong convexity.}
If $f$ is $\mu$-strongly convex and $L$-smooth, gradient descent with $\alpha = 1/L$ satisfies
\handoutdisplay{
f(x_{k+1}) - f^* \le \Big(1 - \frac\mu L\Big)\big(f(x_k) - f^*\big), \qquad f(x_k) - f^* \le \Big(1 - \frac\mu L\Big)^k\big(f(x_0) - f^*\big).
}

\end{theorembox}

\begin{proofbox}
Subtract $f^*$ from both sides of the descent step and bound the gradient term from below by the lemma, $\frac{1}{2L}\|\nabla f(x_k)\|^2 \ge \frac\mu L\big(f(x_k) - f^*\big)$:
\handoutdisplay{
f(x_{k+1}) - f^* \le f(x_k) - f^* - \frac{1}{2L}\|\nabla f(x_k)\|^2 \le \Big(1 - \frac\mu L\Big)\big(f(x_k) - f^*\big).
}
 Applying this $k$ times gives the second inequality.
\end{proofbox}

The factor $1 - \mu/L = 1 - 1/\kappa$ lies in $[0, 1)$. A bound of the form $C q^k$ with $q < 1$ is called a \textbf{linear rate}: on a logarithmic scale the bound falls along a straight line.

\begin{whybox}
Strong convexity keeps the gradient large while the gap is large. With the step $1/L$ every iteration therefore removes at least the fraction $\mu/L$ of the remaining gap.
\end{whybox}

\begin{selfcheckbox}
Take $f(x) = \frac{1}{2}x^\top Q x$ with $\kappa = 1$. What does the theorem give, and what does gradient descent with $\alpha = 1/L$ actually do?
\selfcheckanswer The factor is $1 - \mu/L = 0$, so the gap vanishes after one step. Indeed $\kappa = 1$ means $Q = LI$, and $x_1 = x_0 - \frac1L \cdot L x_0 = 0 = x^*$.
\end{selfcheckbox}

\subsection{The PL condition}\label{sub:rates-pl}

The proof above used strong convexity only once, through the inequality of the lemma. That inequality is the PL condition, and it is enough on its own.

\begin{theorembox}{Theorem: Linear rate under PL.}
Let $f$ be $L$-smooth with $f^* = \inf f > -\infty$, and let it satisfy the Polyak–Łojasiewicz condition with some $\mu > 0$:
\handoutdisplay{
\frac{1}{2}\|\nabla f(x)\|^2 \ge \mu\big(f(x) - f^*\big) \qquad \text{for all } x.
}
 Then gradient descent with $\alpha = 1/L$ satisfies $f(x_k) - f^* \le (1 - \mu/L)^k\big(f(x_0) - f^*\big)$.
\end{theorembox}

\begin{proofbox}
The descent step gives $f(x_{k+1}) \le f(x_k) - \frac{1}{2L}\|\nabla f(x_k)\|^2$, and PL bounds the subtracted term from below by $\frac\mu L\big(f(x_k) - f^*\big)$. Subtracting $f^*$ gives $f(x_{k+1}) - f^* \le (1 - \mu/L)\big(f(x_k) - f^*\big)$, and iterating gives the claim. Convexity was not used anywhere.
\end{proofbox}

PL says that the gradient is large wherever the gap is large. At a stationary point the left-hand side is $0$, so $f(x) = f^*$: \textbf{every stationary point is a global minimizer}. This is usually what convexity is needed for, and PL gives it without convexity.

\begin{examplebox}
$f(x) = x^2 + 3\sin^2 x$ on $\mathbb R$ satisfies PL but is not convex. It is not convex because $f''(x) = 2 + 6\cos 2x$ equals $2 - 6 < 0$ at $x = \pi/2$. It is $L$-smooth with $L = 8$, since $|f''| \le 8$. Its only stationary point is the minimizer $0$: for $0 < x < \tfrac{3}{2}$ both terms of $f'(x) = 2x + 3\sin 2x$ are positive, for $x \ge \tfrac{3}{2}$ we have $2x \ge 3 \ge -3\sin 2x$ with no equality in both places at once, and $f'$ is odd. So the ratio $\frac{1}{2}f'(x)^2/f(x)$ is continuous and positive for $x \ne 0$. It tends to $8$ as $x \to 0$, where $f(x) \approx 4x^2$ and $f'(x) \approx 8x$, and to $2$ as $|x| \to \infty$. Its infimum is therefore positive, and PL holds with $\mu$ equal to it.
\end{examplebox}

\begin{trapbox}
PL and convexity are independent. The example above is PL and not convex. The function $f(x) = x^4$ is convex and not PL: $\frac{1}{2}f'(x)^2 = 8x^6$ is smaller than $\mu x^4$ for every $\mu > 0$ once $|x|$ is small. Near its minimizer $x^4$ is too flat for a linear rate.
\end{trapbox}

\subsection{Condition number and iteration count}\label{sub:rates-kappa}

With $\alpha = 1/L$ the guarantee reads $f(x_k) - f^* \le (1 - 1/\kappa)^k\big(f(x_0) - f^*\big)$. To reduce the gap by a factor $\varepsilon$ it suffices that $(1 - 1/\kappa)^k \le \varepsilon$, that is,
\handoutdisplay{
k \ge \frac{\ln \varepsilon}{\ln(1 - 1/\kappa)},
}
 where both logarithms are negative. Since $\ln(1 - 1/\kappa) \le -1/\kappa$, the right-hand side is at most $\kappa\ln(1/\varepsilon)$, and for large $\kappa$ it is close to it. The count grows linearly in $\kappa$ and only logarithmically in the accuracy: this is the $O(\kappa\log(1/\varepsilon))$ of the lecture.

\begin{examplebox}
Take $\kappa = 1000$ and $\varepsilon = 10^{-6}$. Then $\ln\varepsilon/\ln(1 - 1/\kappa)$ rounds up to $k = 13809$, and the estimate $\kappa\ln(1/\varepsilon)$ gives about $13816$. For a quadratic the optimal step $2/(L + \mu)$ multiplies every eigencomponent of the error by a number of absolute value at most $\rho = (\kappa - 1)/(\kappa + 1)$, so the gap shrinks at least by $\rho^2$ per step, and the same reduction needs
\handoutdisplay{
k \ge \frac{\ln\varepsilon}{2\ln\rho},
}
 that is, $k = 3454$. This is about a quarter of the count at $1/L$, because $\ln\rho \approx -2/\kappa$ for large $\kappa$.
\end{examplebox}

These counts are guarantees, and a run can beat them by far. In A2, with $\kappa = 2$, the guarantee gives 20 iterations and the run needs 10; in A3, with $\kappa = 50$, it gives 684 and the run needs 245. Two effects explain the difference. On a quadratic at $\alpha = 1/L$ the $i$-th eigencomponent of the gap shrinks by $(1 - \lambda_i/L)^2$ per step, so the slowest one shrinks by $(1 - \mu/L)^2$, the square of the guaranteed factor; this alone halves the count. And the stiff component is removed in the first step, since its factor is $1 - L/L = 0$: in A3, $f(x_0) = 25.5$ but $f(x_k) = \frac{1}{2}\,(0.98)^{2k}$ for $k \ge 1$.

With the optimal step $2/(L + \mu) = \tfrac{2}{51}$ both components of A3 shrink by exactly $\tfrac{49}{51}$ in absolute value, the bound $f(x_k) \le \rho^{2k} f(x_0)$ holds with equality, and bound and run agree: 173 iterations in both cases.

\begin{selfcheckbox}
At fixed $\kappa$, what happens to the count when $\varepsilon$ is replaced by $\varepsilon^2$?
\selfcheckanswer It doubles: the count is proportional to $\ln(1/\varepsilon)$, and $\ln(1/\varepsilon^2) = 2\ln(1/\varepsilon)$. Twice as many orders of magnitude of the gap cost twice as many iterations; for large $\kappa$ a condition number twice as large roughly doubles the count as well.
\end{selfcheckbox}

% =====================================================================
\section{Reverse-mode automatic differentiation}\label{sec:ad}

Gradient descent needs $\nabla f$ at every step. \textbf{Automatic differentiation} (AD) computes it from the program that evaluates $f$, exactly up to rounding, by applying the chain rule to one elementary operation at a time. Its reverse mode delivers the whole gradient for a small constant multiple of the cost of one evaluation of $f$, however many variables there are. Everything here is done by hand on small graphs.

\subsection{A computation as a graph}\label{sub:ad-graph}

A program evaluates $f$ through a sequence of intermediate variables, each an elementary function of earlier ones. For the example of the lecture, $z = \sin(xy) + x^2$:
\handoutdisplay{
u = xy, \qquad v = \sin u, \qquad w = x^2, \qquad z = v + w.
}
 The nodes are the inputs $x, y$ and the intermediates $u, v, w, z$; an edge $a \to b$ means that $b$ uses $a$. Every variable uses only earlier ones, so there are no cycles: the graph is a \textbf{directed acyclic graph} (DAG). The order $x, y, u, v, w, z$ in which the program computes the nodes is a \textbf{topological order}: every edge points forward in it.

Each node knows its local partial derivatives: $\partial u/\partial x = y$, $\partial u/\partial y = x$, $\partial v/\partial u = \cos u$, $\partial w/\partial x = 2x$, $\partial z/\partial v = \partial z/\partial w = 1$. AD combines these numbers by the chain rule; it never builds a formula for $\partial z/\partial x$ as a whole.

\handoutfigure{0.95}{fig_graph_c1}{The graph of C1 at $(x, y) = (2,\ -1)$. Forward values: $u = -2$, $v = -0.909297$, $w = 4$, $z = 3.090703$. Adjoints of the reverse pass: $\bar z = \bar v = \bar w = 1$, $\bar u = \cos u = -0.416147$, $\bar x = 4.416147$, $\bar y = -0.832294$.}

\subsection{Forward mode}\label{sub:ad-forward}

Choose a direction $(\dot x, \dot y)$ of the inputs and carry, next to every value, its \textbf{tangent}: the derivative of that value along the chosen direction. The tangents follow the graph forward with the same local derivatives:
\handoutdisplay{
\dot u = y\dot x + x\dot y, \qquad \dot v = \cos(u)\,\dot u, \qquad \dot w = 2x\,\dot x, \qquad \dot z = \dot v + \dot w.
}


For a map $F: \mathbb R^n \to \mathbb R^m$ with Jacobian $J$, one forward sweep with the seed $\dot x \in \mathbb R^n$ returns the \textbf{Jacobian–vector product} $J\dot x$. For a scalar $f$ this is the directional derivative $\nabla f(x)^\top \dot x$, a single number, so the full gradient needs $n$ sweeps, one per basis vector $e_1, \dots, e_n$.

\begin{examplebox}
At the point of C1 with the seed $\dot x = 1$, $\dot y = 0$ one sweep gives $\dot z = y\cos u + 2x = 4.416147$, the partial derivative in $x$. The derivative in $y$ needs a second sweep with $\dot x = 0$, $\dot y = 1$.
\end{examplebox}

\subsection{Reverse mode}\label{sub:ad-reverse}

For a scalar output $z$, the \textbf{adjoint} of a node $v$ is $\bar v = \partial z/\partial v$, the derivative of the output with respect to that node through everything downstream of it. Start with $\bar z = 1$ and every other adjoint equal to $0$, visit the nodes in reverse topological order, and for every edge $u \to v$ add the contribution
\handoutdisplay{
\bar u \mathrel{+}= \bar v\, \frac{\partial v}{\partial u}.
}
 When the sweep reaches a node, all nodes that use it have been processed, so its adjoint is complete. In one formula,
\handoutdisplay{
\bar v_i = \sum_{j:\ i \to j} \bar v_j\, \frac{\partial v_j}{\partial v_i},
}
 the sum over all nodes $j$ that use $v_i$.

In the example: $\bar z = 1$; the node $z = v + w$ gives $\bar v = 1$ and $\bar w = 1$; $w = x^2$ gives $\bar x \mathrel{+}= 2x$; $v = \sin u$ gives $\bar u = \cos u$; $u = xy$ gives $\bar x \mathrel{+}= y\cos u$ and $\bar y \mathrel{+}= x\cos u$. Hence
\handoutdisplay{
\bar x = y\cos u + 2x = \frac{\partial z}{\partial x}, \qquad \bar y = x\cos u = \frac{\partial z}{\partial y}.
}
 One backward sweep has produced both partial derivatives.

For $F: \mathbb R^n \to \mathbb R^m$ a reverse sweep started from a seed $\bar z \in \mathbb R^m$ returns the \textbf{vector–Jacobian product} $\bar z^\top J$. For a scalar loss, $m = 1$ and $\bar z = 1$, this is the whole gradient in one sweep, whatever $n$ is. Reverse mode applied to a model with a scalar loss and many parameters is \textbf{backpropagation}.

\subsection{Why adjoints accumulate}\label{sub:ad-accumulate}

A node that feeds several others receives a contribution along each outgoing edge, and the multivariable chain rule adds them. For $z = x^2 + x^3$ with the nodes $a = x^2$, $b = x^3$, $z = a + b$:
\handoutdisplay{
\frac{dz}{dx} = \frac{\partial z}{\partial a}\frac{da}{dx} + \frac{\partial z}{\partial b}\frac{db}{dx} = 2x + 3x^2.
}
 This is why the rule is written with $\mathrel{+}=$: the adjoint $\bar x$ collects $2x$ from $a$ and $3x^2$ from $b$. An implementation that assigns instead of adding keeps only the contribution of the edge it visits last, $2x$ or $3x^2$, and the wrong answer even depends on the visit order.

In general $\partial z/\partial x$ is the sum, over all paths from $x$ to $z$, of the products of the local derivatives along the path. The reverse sweep computes this sum without listing the paths: the adjoint of each node already holds the sum over all paths from that node to $z$. In C1 the input $x$ feeds $u$ and $w$, so $\bar x$ has two terms; in C2 it feeds both $xy$ and $xy + x$.

\begin{trapbox}
Frameworks accumulate parameter gradients across calls for the same reason. Before a new gradient is computed, the stored adjoints must be reset to zero, or the new gradient is added to the old one.
\end{trapbox}

\subsection{Cost and memory}\label{sub:ad-cost}

Let $f: \mathbb R^n \to \mathbb R$ be computed by $N$ elementary operations. A forward or a reverse sweep does a bounded amount of work per operation, so each costs a small constant multiple of one evaluation of $f$. The difference between the modes is the number of sweeps.

\begin{center}\small
\begin{tabular}{@{}>{\raggedright\arraybackslash\hspace{0pt}}p{0.27\linewidth}>{\raggedright\arraybackslash\hspace{0pt}}p{0.3\linewidth}>{\raggedright\arraybackslash\hspace{0pt}}p{0.32\linewidth}@{}}
\toprule
Property & Forward mode & Reverse mode \\
\midrule
Seed & a direction $\dot x \in \mathbb R^n$ & an output weight $\bar z \in \mathbb R^m$ \\
One sweep gives & $J\dot x$, a Jacobian–vector product & $\bar z^\top J$, a vector–Jacobian product \\
Gradient of $f: \mathbb R^n \to \mathbb R$ & $n$ sweeps & one forward and one reverse sweep \\
Jacobian of $F: \mathbb R^n \to \mathbb R^m$ & $n$ sweeps & $m$ sweeps \\
Extra memory & one tangent per value in use & every intermediate value (the tape) \\
\bottomrule
\end{tabular}
\end{center}

Reverse mode is cheap in operations but not in memory. The local derivatives need the forward values ($\partial v/\partial u = \cos u$ needs $u$), so the forward sweep must record every intermediate value, the \textbf{tape}, before the reverse sweep can start, and the memory grows with the length of the computation.

\textbf{Checkpointing} trades computation for memory. The forward sweep stores only some intermediate values, the checkpoints; during the reverse sweep each segment between two checkpoints is recomputed from the earlier one when its values are needed. With about $\sqrt N$ checkpoints spaced $\sqrt N$ operations apart, the memory drops to the order of $\sqrt N$ stored values at the price of roughly one extra forward sweep.

\subsection{Backpropagation for the logistic loss}\label{sub:ad-logistic}

One observation has features $x \in \mathbb R^d$ and a label $y \in \{0, 1\}$. The model computes
\handoutdisplay{
s = w^\top x + b, \qquad p = \sigma(s) = \frac{1}{1 + e^{-s}}, \qquad \ell = -y\log p - (1 - y)\log(1 - p).
}
 Since $\log p = s - \log(1 + e^s)$ and $\log(1 - p) = -\log(1 + e^s)$, the loss is $\ell = \log(1 + e^s) - ys$, the form used in C3.

The reverse pass starts with $\bar\ell = 1$. The node $\ell(s)$ gives
\handoutdisplay{
\bar s = \frac{d\ell}{ds} = \frac{e^s}{1 + e^s} - y = \sigma(s) - y,
}
 and the node $s = w^\top x + b$ passes it on to the parameters: $\bar w = \bar s\, x$ and $\bar b = \bar s$. So
\handoutdisplay{
\nabla_w \ell = \big(\sigma(s) - y\big)\,x, \qquad \frac{\partial\ell}{\partial b} = \sigma(s) - y,
}
 and the forward value $p = \sigma(s)$ is reused. For the mean loss $\frac1n\sum_i \ell_i$ the contributions of the observations are added and divided by $n$: $\nabla_w \frac1n\sum_i \ell_i = \frac1n\sum_i \big(\sigma(s_i) - y_i\big)\,x_i$.

With labels $\tilde y = 2y - 1 \in \{-1, +1\}$ the same loss reads
\handoutdisplay{
\ell = \log\big(1 + e^{-\tilde y s}\big).
}
 The two cases check it: $y = 1$, $\tilde y = 1$ gives $\log(1 + e^{-s}) = \log(1 + e^s) - s$, and $y = 0$, $\tilde y = -1$ gives $\log(1 + e^s)$. Its derivative $d\ell/ds = -\tilde y\,\sigma(-\tilde y s)$ equals $\sigma(s) - y$ in both cases. The two conventions describe one function.

\begin{trapbox}
$\sigma(s) - y$ is the derivative with respect to $s$, not with respect to $w$. The gradient in $w$ carries the factor $x$; without it the result is a number where a vector in $\mathbb R^d$ is expected.
\end{trapbox}

\subsection{Symbolic, numerical and automatic differentiation}\label{sub:ad-compare}

\begin{center}\small
\begin{tabular}{@{}>{\raggedright\arraybackslash\hspace{0pt}}p{0.15\linewidth}>{\raggedright\arraybackslash\hspace{0pt}}p{0.14\linewidth}>{\raggedright\arraybackslash\hspace{0pt}}p{0.24\linewidth}>{\raggedright\arraybackslash\hspace{0pt}}p{0.19\linewidth}>{\raggedright\arraybackslash\hspace{0pt}}p{0.13\linewidth}@{}}
\toprule
Method & Core idea & Accuracy & Typical issue & Cost of a gradient in $\mathbb R^n$ \\
\midrule
Symbolic & manipulates the formula into a formula for the derivative & exact & expression swell: the derivative formula can be far larger than the program & evaluating that formula, whose size is not controlled \\
Finite differences & perturbs the inputs and divides differences of values & approximate, at best about $\varepsilon^{2/3}$ relative for central differences, with $\varepsilon$ the machine epsilon & truncation for large steps, cancellation for small ones & $2n$ evaluations of $f$ for central differences \\
Automatic, reverse mode & applies the chain rule to the program, node by node & exact up to rounding & memory for the tape; control flow and non-smooth operations & a small constant multiple of one evaluation of $f$ \\
\bottomrule
\end{tabular}
\end{center}

The central difference
\handoutdisplay{
\frac{\partial f}{\partial x_i}(x) \approx \frac{f(x + he_i) - f(x - he_i)}{2h}
}
 has two errors. Taylor's formula $f(x \pm he_i) = f \pm h\,\partial_i f + \frac{h^2}{2}\partial_i^2 f \pm \frac{h^3}{6}\partial_i^3 f + O(h^4)$ makes the even terms cancel, so the \textbf{truncation} error is $\frac{h^2}{6}|\partial_i^3 f| + O(h^3)$, of order $h^2$. Each computed value of $f$ carries a \textbf{round-off} error of about $\varepsilon|f|$, where $\varepsilon$ is the machine epsilon of double precision. The difference of two nearly equal numbers keeps these absolute errors, and dividing by $2h$ gives an error of order $\varepsilon|f|/h$. The total
\handoutdisplay{
E(h) \approx \frac{h^2}{6}M + \frac{\varepsilon|f|}{h}, \qquad M = |\partial_i^3 f|,
}
 is smallest where $E'(h) = \frac{hM}{3} - \frac{\varepsilon|f|}{h^2} = 0$, at $h^3 = 3\varepsilon|f|/M$. So the best step is of order $\varepsilon^{1/3} \approx 6 \cdot 10^{-6}$ and the best error of order $\varepsilon^{2/3}$. The one-sided difference $\big(f(x + he_i) - f(x)\big)/h$ has truncation error of order $h$, and the same argument gives the best step of order $\sqrt\varepsilon$.

Demo D3 measures this for $\partial z/\partial x$ of C1 at $(x, y) = (2,\ -1)$. The relative error falls like $h^2$ from $1.6 \cdot 10^{-4}$ at $h = 10^{-1}$, reaches its minimum $5.7 \cdot 10^{-12}$ at $h = 10^{-5}$, close to $\varepsilon^{1/3}$, and grows again to $1.4 \cdot 10^{-7}$ at $h = 10^{-10}$ as round-off takes over.

\handoutfigure{0.72}{fig_fd_error}{Relative error of the central difference for $\partial z/\partial x$ of C1 at $(x, y) = (2,\ -1)$, for $h$ from $10^{-1}$ to $10^{-10}$. For large $h$ the truncation error dominates and falls as $h^2$; for small $h$ round-off dominates and grows as $1/h$.}

AD has no step to choose and no truncation error: it evaluates the exact derivative of the program in floating point. Finite differences remain useful as an independent check of a gradient, at a few points, with $h$ near $\varepsilon^{1/3}$ and a relative error measure.

\subsection{Problems of block C}\label{sub:ad-problems}

The three seminar problems on the reverse pass, with complete solutions.

\problemhead{C1}{The lecture's graph at a point}{On the board}
\textbf{Statement.} Take $z = \sin(xy) + x^2$ at $(x, y) = (2,\ -1)$ with the intermediate nodes $u = xy$, $w = x^2$, $v = \sin u$. List the nodes in topological order, run the forward pass, then the reverse pass from $\bar z = 1$. Report $z$, $\partial z/\partial x$ and $\partial z/\partial y$, and check both derivatives analytically.

\textbf{Solution.}
\begingroup\sloppy
\begin{enumerate}[label=\textbf{\arabic*.}, leftmargin=*]
  \item Topological order, every node after its inputs: $x$, $y$, $u = xy$, $w = x^2$, $v = \sin u$, $z = v + w$. The input $x$ has two outgoing edges, into $u$ and into $w$.
  \item Forward pass at $(x, y) = (2,\ -1)$, keeping every value:
\handoutdisplay{
u = -2,\quad w = 4,\quad v = \sin(-2) = -0.909297,\quad z = v + w = 3.090703.
}
 The reverse pass needs $u$, $x$ and $y$ again.
  \item Reverse pass, node $z = v + w$: seed $\bar z = 1$. Both local derivatives are $1$, so $\bar v = 1$ and $\bar w = 1$.
  \item Node $v = \sin u$: with $u$ in radians, $\bar u \mathrel{+}= \bar v \cos u = \cos(-2) = -0.416147$.
  \item Node $w = x^2$: the local derivative is $2x$, so $\bar x \mathrel{+}= \bar w \cdot 2x = 2 \cdot 2$. This is the first of the two contributions to $\bar x$.
  \item Node $u = xy$: the local derivatives are $y$ and $x$, so
\handoutdisplay{
\bar x \mathrel{+}= \bar u\,y = (-0.416147)(-1) = 0.416147, \qquad \bar y \mathrel{+}= \bar u\,x = (-0.416147) \cdot 2 = -0.832294.
}
 With both contributions $\bar x = 0.416147 + 2 \cdot 2 = 4.416147$; overwriting would leave only $0.416147$.
  \item Analytic check: $\partial z/\partial x = y\cos(xy) + 2x$ and $\partial z/\partial y = x\cos(xy)$ give the same numbers. The reverse pass evaluated these formulas term by term: $\bar u\,y$ is $y\cos(xy)$ and $\bar w \cdot 2x$ is $2x$. One sweep produced both partial derivatives.
\end{enumerate}
\endgroup
\begin{answerbox}
$z = 3.090703$, $\partial z/\partial x = 4.416147$, $\partial z/\partial y = -0.832294$.
\end{answerbox}

\problemhead{C2}{An input used twice}{On your own}
\textbf{Statement.} Take $q = (xy + x)^2$ at $(x, y) = (2,\ 3)$ with the nodes $a = xy$, $b = a + x$, $q = b^2$. Compute $q$ by a forward pass and $\partial q/\partial x$, $\partial q/\partial y$ by one reverse pass, with pen and paper.

\textbf{Solution.}
\begingroup\sloppy
\begin{enumerate}[label=\textbf{\arabic*.}, leftmargin=*]
  \item Topological order: $x$, $y$, $a = xy$, $b = a + x$, $q = b^2$. Two edges leave $x$: one into $a$ and one directly into $b$.
  \item Forward pass at $(x, y) = (2,\ 3)$: $a = xy = 6$, $b = a + x = 8$, $q = b^2 = 64$.
  \item Reverse pass: $\bar q = 1$, and the node $q = b^2$ gives $\bar b = \bar q \cdot 2b = 16$.
  \item Node $b = a + x$: both local derivatives are $1$, so $\bar a \mathrel{+}= \bar b = 16$ and $\bar x \mathrel{+}= \bar b = 16$.
  \item Node $a = xy$: $\bar x \mathrel{+}= \bar a\,y = 48$ and $\bar y \mathrel{+}= \bar a\,x = 32$. So $\bar x = 16 + 48 = 64$; with \texttt{=} in place of \texttt{+=} it would be $48$, the last contribution only.
  \item Analytic check: $q = x^2(y + 1)^2$, so $\partial q/\partial x = 2x(y + 1)^2 = 64$ and $\partial q/\partial y = 2x^2(y + 1) = 32$. The value $q$ and $\partial q/\partial x$ agree here only because $x = 2$: their ratio is $x/2$.
\end{enumerate}
\endgroup
\begin{answerbox}
$q = 64$, $\partial q/\partial x = 48 + 16 = 64$, $\partial q/\partial y = 32$.
\end{answerbox}

\problemhead{C3}{The logistic loss node}{On your own}
\textbf{Statement.} For one example with score $s$ and label $y \in \{0, 1\}$ the logistic loss is $\ell(s) = \log(1 + e^s) - ys$, and $\sigma(s) = 1/(1 + e^{-s})$. At $s = 0.5$ and $y = 1$, split $\ell$ into elementary nodes, run the forward and the reverse pass, and report $\ell$, $\sigma(s)$ and $d\ell/ds$. Then rewrite the loss for labels $\tilde y = 2y - 1 \in \{-1, 1\}$ as a function of $\tilde y s$ alone, and check it for both labels.

\textbf{Solution.}
\begingroup\sloppy
\begin{enumerate}[label=\textbf{\arabic*.}, leftmargin=*]
  \item Nodes in topological order: $s$, $u = e^s$, $v = 1 + u$, $w = \log v$, $\ell = w - ys$. The label $y$ is a constant. The score $s$ feeds two nodes, $u$ and $\ell$, so $\bar s$ will have two terms.
  \item Forward pass at $s = 0.5$, $y = 1$:
\handoutdisplay{
\begin{gathered}
u = e^s = 1.648721,\quad v = 1 + 1.648721,\quad w = \log v = 0.974077, \\[6pt]
\ell = w - ys = 0.974077 - 1 \cdot 0.5 = 0.474077.
\end{gathered}
}

  \item Reverse pass: $\bar\ell = 1$. The node $\ell = w - ys$ has local derivative $1$ in $w$ and $-y$ in $s$, so $\bar w = 1$ and $\bar s \mathrel{+}= -y = -1$.
  \item Through the logarithm: $\bar v = \bar w / v = \dfrac{1}{1 + e^s}$, then $\bar u = \bar v$, then
\handoutdisplay{
\bar s \mathrel{+}= \bar u\,e^s = \frac{e^s}{1 + e^s} = \frac{1}{1 + e^{-s}} = \sigma(s) = 0.622459.
}

  \item Both terms together: $d\ell/ds = \bar s = \sigma(s) - y = 0.622459 - 1 = -0.377541$. This is the lecture's $p - y$ with $p = \sigma(s)$. It is negative, so a gradient step raises the score of a positive example.
  \item Labels $\pm 1$: put $\tilde y = 2y - 1$. For $y = 1$, $\tilde y = 1$ and $\log(1 + e^s) - s = \log\dfrac{1 + e^s}{e^s} = \log(1 + e^{-s})$. For $y = 0$, $\tilde y = -1$ and the loss is $\log(1 + e^s)$ unchanged. Both cases read
\handoutdisplay{
\ell = \log\big(1 + e^{-\tilde y s}\big).
}

  \item Check at the point: $\tilde y = 1$ and $\log(1 + e^{-\tilde y s}) = 0.474077$, the same number as $\ell$. The two conventions are one function, so they share every value and the derivative $\sigma(s) - y$. The lecture writes the loss with labels in $\{0, 1\}$, and the $\pm 1$ form is the other common convention.
\end{enumerate}
\endgroup
\begin{answerbox}
$\ell = 0.474077$, $\sigma(s) = 0.622459$, $d\ell/ds = \sigma(s) - y = -0.377541$; with $\tilde y = 2y - 1$ the loss is $\log(1 + e^{-\tilde y s})$.
\end{answerbox}

% =====================================================================
\section{Colloquium}\label{sec:colloquium}

At the colloquium definitions come first, so the checklist of definitions from the first section comes first here too. Then come the four proofs with their key inequalities and three ticket-style problems with full solutions.

\subsection{Definitions checklist}\label{sub:colloquium-definitions}

For each definition give the name, the formula and the quantifiers: for which points, or pairs of points, the formula must hold.

\begin{itemize}
  \item Local minimizer: $x^*$ such that there is $\varepsilon > 0$ with $f(x) \ge f(x^*)$ for all $x$ with $\|x - x^*\| < \varepsilon$.
  \item Global minimizer: $x^*$ such that $f(x) \ge f(x^*)$ for all $x$.
  \item Stationary point: a point $x$ at which $f$ is differentiable and $\nabla f(x) = 0$.
  \item Convex set: $(1 - t)x + ty \in C$ for all $x, y \in C$ and $t \in [0, 1]$.
  \item Convex function (Jensen): $f\big((1 - t)x + ty\big) \le (1 - t)f(x) + t f(y)$ for all $x, y$ and $t \in [0, 1]$.
  \item Strictly convex function: the same inequality is strict for all $x \ne y$ and $t \in (0, 1)$.
  \item $\mu$-strongly convex function, $\mu > 0$: $f(y) \ge f(x) + \nabla f(x)^\top(y - x) + \frac\mu2\|y - x\|^2$ for all $x, y$; equivalently, $f - \frac\mu2\|\cdot\|^2$ is convex.
  \item $L$-smooth function: $\|\nabla f(x) - \nabla f(y)\| \le L\|x - y\|$ for all $x, y$.
  \item Condition number: $\kappa = L/\mu \ge 1$; for $f = \frac{1}{2}x^\top Q x$ it is $\lambda_{\max}(Q)/\lambda_{\min}(Q)$.
  \item Descent direction at $x$: $d$ with $\nabla f(x)^\top d < 0$.
  \item PL condition with $\mu > 0$: $\frac{1}{2}\|\nabla f(x)\|^2 \ge \mu\big(f(x) - f^*\big)$ for all $x$.
\end{itemize}

\subsection{Four proofs}\label{sub:colloquium-proofs}

\begin{notebox}{Descent lemma}
If $\nabla f$ is $L$-Lipschitz, then $f(y) \le f(x) + \nabla f(x)^\top(y - x) + \frac L2\|y - x\|^2$. Put $r(t) = f(x + t(y - x))$, write $f(y) - f(x) = \int_0^1 r'(t)\,dt$ with $r'(t) = \nabla f(x + t(y - x))^\top(y - x)$, subtract $\nabla f(x)^\top(y - x)$ and bound the rest by Cauchy–Schwarz and the Lipschitz condition:
\handoutdisplay{
f(y) - f(x) - \nabla f(x)^\top(y - x) = \int_0^1 \big(\nabla f(x + t(y - x)) - \nabla f(x)\big)^\top (y - x)\,dt \le \int_0^1 Lt\|y - x\|^2\,dt = \frac L2\|y - x\|^2.
}

\end{notebox}

\begin{notebox}{Armijo termination}
If $\nabla f$ is $L$-Lipschitz and $g = \nabla f(x) \ne 0$, backtracking along $d = -g$ stops after finitely many reductions. The descent lemma at $y = x - \alpha g$ gives
\handoutdisplay{
f(x - \alpha g) \le f(x) - \alpha\Big(1 - \frac{L\alpha}{2}\Big)\|g\|^2,
}
 and the right-hand side is at most $f(x) - c\alpha\|g\|^2$ as soon as $1 - L\alpha/2 \ge c$, that is, for every $\alpha \le 2(1 - c)/L$. The trials $\beta^j\alpha_0$ tend to $0$, so one of them enters this interval. The accepted step is either $\alpha_0$ or $\beta$ times a rejected trial, and a rejected trial is larger than $2(1 - c)/L$, so $\alpha \ge \min\{\alpha_0,\ 2\beta(1 - c)/L\}$.
\end{notebox}

\begin{notebox}{Rate for convex functions}
For convex $L$-smooth $f$ and $\alpha = 1/L$, $f(x_k) - f^* \le L\|x_0 - x^*\|^2/(2k)$. Combine the descent step $f(x_{k+1}) \le f(x_k) - \frac{1}{2L}\|g_k\|^2$, convexity $f(x_k) - f^* \le g_k^\top(x_k - x^*)$ and the three-point identity for $x_{k+1} = x_k - \frac1L g_k$ into the key inequality
\handoutdisplay{
f(x_{k+1}) - f^* \le \frac L2\Big(\|x_k - x^*\|^2 - \|x_{k+1} - x^*\|^2\Big).
}
 Sum it over the first $k$ steps, where the right-hand side telescopes, and use that $f(x_j)$ does not increase.
\end{notebox}

\begin{notebox}{Linear rate under strong convexity and under PL}
Strong convexity implies PL: minimize the lower bound $f(x) + \nabla f(x)^\top(y - x) + \frac\mu2\|y - x\|^2$ over $y$ to get $f^* \ge f(x) - \frac{1}{2\mu}\|\nabla f(x)\|^2$. Under PL and $\alpha = 1/L$ the key inequality is
\handoutdisplay{
f(x_{k+1}) - f^* \le f(x_k) - f^* - \frac{1}{2L}\|\nabla f(x_k)\|^2 \le \Big(1 - \frac\mu L\Big)\big(f(x_k) - f^*\big),
}
 and iterating it gives $(1 - \mu/L)^k$. The second proof uses only PL, not convexity.
\end{notebox}

\subsection{Ticket T1 on conditioning and rates}\label{sub:colloquium-t1}

\begin{selfcheckbox}
Let $f(x) = \frac{1}{2}x^\top Q x$ with $Q = \operatorname{diag}(4,\ 100)$. Find $\mu$, $L$ and $\kappa$. For which constant steps does gradient descent converge from every starting point? How many iterations does the guarantee at $\alpha = 1/L$ need to reduce $f(x_k) - f^*$ by the factor $10^{-6}$? Find the optimal constant step, its contraction factor and the number of iterations it needs for the same reduction.
\selfcheckanswer $\mu = \lambda_{\min}(Q) = 4$, $L = \lambda_{\max}(Q) = 100$, $\kappa = L/\mu = 25$. In eigen-coordinates the iteration is $x_{k+1,i} = (1 - \alpha\lambda_i)x_{k,i}$, so it converges from every start exactly when $|1 - \alpha\lambda_i| < 1$ for both $i$, that is, for $0 < \alpha < 2/L = 0.02$. At $\alpha = 1/L = 0.01$ the guarantee $f(x_k) - f^* \le (1 - \mu/L)^k\big(f(x_0) - f^*\big)$ has the factor $1 - \mu/L = 0.96$, and $(0.96)^k \le 10^{-6}$ requires $k \ge \ln(10^{-6})/\ln(0.96)$, so $k = 339$. The optimal constant step minimizes the worst factor $\max\{|1 - \alpha\mu|,\ |1 - \alpha L|\}$. The first term decreases in $\alpha$ and the second increases once $\alpha > 1/L$, so the minimum is where $1 - \alpha\mu = \alpha L - 1$:
\handoutdisplay{
\alpha^* = \frac{2}{L + \mu} = \tfrac{1}{52} \approx 0.019231, \qquad \rho = \frac{L - \mu}{L + \mu} = \frac{\kappa - 1}{\kappa + 1} = \tfrac{12}{13} \approx 0.923077.
}
 Both components of $x_k$ shrink by exactly $\rho$ in absolute value, so $f(x_k) = \rho^{2k} f(x_0)$, and $\rho^{2k} \le 10^{-6}$ requires $k \ge \ln(10^{-6})/(2\ln\rho)$, so $k = 87$: almost four times fewer iterations than the guarantee at $1/L$.
\end{selfcheckbox}

\subsection{Ticket T2 on Armijo backtracking}\label{sub:colloquium-t2}

\begin{selfcheckbox}
Run Armijo backtracking on $f(x) = \frac{1}{2}x^\top \operatorname{diag}(2,\ 20)\,x$ from $x_0 = (1,\ 1)$ along $d = -\nabla f(x_0)$ with $\alpha_0 = 1$, $\beta = 0.5$ and $c = 0.25$. List the trials, find the accepted step and the number of halvings, and check the result against the termination theorem.
\selfcheckanswer Here $L = \lambda_{\max} = 20$, $f(x_0) = 11$, $g = \nabla f(x_0) = (2,\ 20)$ and $\|g\|^2 = 404$. A trial $\alpha$ is accepted when $f(x_0 - \alpha g) \le f(x_0) - c\alpha\|g\|^2$:
\handoutdisplay{
\begin{array}{c|c|c|c} \alpha & f(x_0 - \alpha g) & f(x_0) - c\alpha\|g\|^2 & \text{accepted} \\ \hline 1 & 3611 & -90 & \text{no} \\ 0.5 & 810 & -39.5 & \text{no} \\ 0.25 & 160.25 & -14.25 & \text{no} \\ 0.125 & 23.0625 & -1.625 & \text{no} \\ 0.0625 & 1.390625 & 4.6875 & \text{yes} \end{array}
}
 The loop stops after 4 halvings with $\alpha = 0.0625$. The theorem guarantees acceptance for every $\alpha \le 2(1 - c)/L = 0.075$, so the loop cannot go past the first trial below this value, which is $0.0625$. Since at least one trial was rejected, the accepted step also satisfies $\alpha \ge \beta \cdot 2(1 - c)/L = 0.0375$. Both checks agree with the run: $0.0375 \le 0.0625 \le 0.075$.
\end{selfcheckbox}

\subsection{Ticket T3 on a reverse pass}\label{sub:colloquium-t3}

\begin{selfcheckbox}
Compute $f(x, y) = x e^y + y^2$ and both partial derivatives at $(x, y) = (2,\ 1)$ by a reverse pass: list the nodes in topological order, compute the forward values, then the adjoints, and check the result against the analytic gradient.
\selfcheckanswer Nodes in topological order: the inputs $x, y$, then $a = e^y$, $b = xa$, $c = y^2$, $f = b + c$. Forward values: $a = e \approx 2.718282$, $b = 5.436564$, $c = 1$, $f = 6.436564$. Reverse pass from $\bar f = 1$: the node $f = b + c$ gives $\bar b = 1$ and $\bar c = 1$; $c = y^2$ gives $\bar y \mathrel{+}= \bar c \cdot 2y = 2$; $b = xa$ gives $\bar x \mathrel{+}= \bar b\, a = 2.718282$ and $\bar a \mathrel{+}= \bar b\, x = 2$; $a = e^y$ gives $\bar y \mathrel{+}= \bar a\, e^y = 5.436564$. The input $y$ feeds two nodes, so its adjoint accumulates: $\bar y = 2 + 5.436564 = 7.436564$. Analytically $\partial f/\partial x = e^y = e$ and $\partial f/\partial y = x e^y + 2y = 2e + 2$, which at this point are $2.718282$ and $7.436564$. Overwriting $\bar y$ instead of adding would give $2$ or $5.436564$, depending on the order.
\end{selfcheckbox}

% =====================================================================
\section{Useful constants}\label{sec:constants}

The step $1/L$, the stability boundary $2/L$ and the rate $1 - \mu/L$ all need the constants of the problem. For a quadratic they are eigenvalues; for logistic regression $L$ has a simple upper bound.

\subsection{Logistic regression}\label{sub:constants-logistic}

The data are $x_1, \dots, x_n \in \mathbb R^d$, the rows of $X \in \mathbb R^{n \times d}$, with labels $\tilde y_i \in \{-1, +1\}$. The mean logistic loss is
\handoutdisplay{
f(w) = \frac1n\sum_{i=1}^n \log\big(1 + e^{-\tilde y_i w^\top x_i}\big),
}
 or, in the $\{0, 1\}$ form, $\frac1n\sum_i\big(\log(1 + e^{s_i}) - y_i s_i\big)$ with $s_i = w^\top x_i$ and $\tilde y_i = 2y_i - 1$; the two are the same function. A bias is absorbed into $w$ by a constant feature $1$ in every $x_i$.

\begin{theorembox}{Proposition: Smoothness of the logistic loss.}
With $\sigma_i = \sigma(w^\top x_i)$,
\handoutdisplay{
\nabla^2 f(w) = \frac1n X^\top D X, \qquad D = \operatorname{diag}\big(\sigma_i(1 - \sigma_i)\big),
}
 and $0 \preceq \nabla^2 f(w) \preceq \frac{1}{4n}X^\top X$ for every $w$. Hence $f$ is convex and $L$-smooth with
\handoutdisplay{
L \le \frac{\lambda_{\max}(X^\top X)}{4n}.
}

\end{theorembox}

\begin{proofbox}
In the $\{0, 1\}$ form the $i$-th term is $\varphi(s_i) - y_i s_i$ with $\varphi(s) = \log(1 + e^s)$, $\varphi'(s) = \sigma(s)$ and $\varphi''(s) = \sigma(s)\big(1 - \sigma(s)\big)$; the linear term has no curvature. By the chain rule with $s_i = w^\top x_i$ the Hessian of the $i$-th term is $\sigma_i(1 - \sigma_i)\,x_i x_i^\top$, and the mean of these matrices is $\frac1n X^\top D X$. Since $p(1 - p) \le \frac{1}{4}$ for $p \in [0, 1]$, with equality at $p = \frac{1}{2}$, for every $v$
\handoutdisplay{
0 \le v^\top \nabla^2 f(w)\, v = \frac1n\sum_i \sigma_i(1 - \sigma_i)(x_i^\top v)^2 \le \frac{1}{4n}\sum_i (x_i^\top v)^2 = \frac{1}{4n}\|Xv\|^2 \le \frac{\lambda_{\max}(X^\top X)}{4n}\|v\|^2.
}
 A twice differentiable function with $0 \preceq \nabla^2 f \preceq L I$ everywhere is convex, and its gradient is $L$-Lipschitz because $\nabla f(y) - \nabla f(x) = \int_0^1 \nabla^2 f\big(x + t(y - x)\big)(y - x)\,dt$.
\end{proofbox}

Here $\lambda_{\max}(X^\top X) = \|X\|_2^2$, the square of the largest singular value of $X$. The bound is attained at $w = 0$, where every $\sigma_i = \frac{1}{2}$ and $\nabla^2 f(0) = \frac{1}{4n}X^\top X$, so no smaller constant is valid on the whole space. The loss is not strongly convex: along $w = tv$ with $x_i^\top v \ne 0$ for all $i$ the weights $\sigma_i(1 - \sigma_i)$ tend to $0$ as $t \to \infty$, and so does the Hessian.

With a penalty, $f_\lambda(w) = f(w) + \lambda\|w\|^2$ with a weight $\lambda > 0$ (a number, not an eigenvalue), the Hessian gains $\nabla^2\big(\lambda\|w\|^2\big) = 2\lambda I$:
\handoutdisplay{
2\lambda I \preceq \nabla^2 f_\lambda(w) \preceq \Big(\frac{\lambda_{\max}(X^\top X)}{4n} + 2\lambda\Big) I.
}
 So $f_\lambda$ is $2\lambda$-strongly convex and $L$-smooth with $L \le \frac{\lambda_{\max}(X^\top X)}{4n} + 2\lambda$, and
\handoutdisplay{
\kappa = \frac L\mu \le \frac{\lambda_{\max}(X^\top X)/(4n) + 2\lambda}{2\lambda} = 1 + \frac{\lambda_{\max}(X^\top X)}{8n\lambda}.
}
 Gradient descent with $\alpha = 1/L$ then converges linearly, and a smaller $\lambda$ means a larger bound on $\kappa$ and a slower guaranteed rate.

\begin{trapbox}
Some texts write the penalty as $\frac\lambda2\|w\|^2$. Then the Hessian gains $\lambda I$, not $2\lambda I$, the problem is $\lambda$-strongly convex, and every constant above changes accordingly.
\end{trapbox}

\subsection{The quadratics of the seminar}\label{sub:constants-quadratics}

For $f(x) = \frac{1}{2}x^\top Q x$ with diagonal $Q$: $\mu = \lambda_{\min}(Q)$, $L = \lambda_{\max}(Q)$, gradient descent converges for $0 < \alpha < 2/L$, the linear guarantee holds at $\alpha = 1/L$, and the optimal constant step is $2/(L + \mu)$.

\begin{center}\small
\begin{tabular}{@{}llrrrrrr@{}}
\toprule
Problem & $Q$ & $\mu$ & $L$ & $\kappa$ & $1/L$ & $2/L$ & $2/(L + \mu)$ \\
\midrule
A1 & $\operatorname{diag}(1,\ 10)$ & $1$ & $10$ & $10$ & $0.1$ & $0.2$ & $\tfrac{2}{11}$ \\
A2 & $\operatorname{diag}(1,\ 2)$ & $1$ & $2$ & $2$ & $0.5$ & $1$ & $\tfrac{2}{3}$ \\
A3 & $\operatorname{diag}(1,\ 50)$ & $1$ & $50$ & $50$ & $0.02$ & $0.04$ & $\tfrac{2}{51}$ \\
\bottomrule
\end{tabular}
\end{center}

B1 uses the matrix of A1; B2 and B3 use the matrix of A3.

\end{document}
